% =====================================================================
%  U06 —— 第 4 章：oper 与一般李代数的中心
%  底本：E. Frenkel, Langlands Correspondence for Loop Groups (2007)
%  PDF p115–p139 = 印刷 p101–p125
% =====================================================================

\chapter{oper 与一般李代数的中心}

现在我们要把定理 3.5.2 推广到任意单李代数 $\mathfrak{g}$ 的情形；该定理描述的是 $\mathfrak{g}=\mathfrak{sl}_2$ 时顶点代数 $V_{\kappa_c}(\mathfrak{g})$ 的中心 $\mathfrak{z}(\hat{\mathfrak{g}})$。第一步，需要把射影联络的概念推广到任意 $\mathfrak{g}$：我们已经看到，在 $\mathfrak{g}=\mathfrak{sl}_2$ 的情形，正是射影联络给出了中心。在 §4.1 中我们重新审视射影联络的概念，用带某种附加结构的平坦 $PGL_2$-丛重新表述它。这使我们能把射影联络推广到任意 $\mathfrak{g}$，推广的形式就是 oper。在 §4.2 中我们将详细讨论 opers 的各种定义与实现，以及坐标变换对它们的作用。

然后我们在 §4.3 中陈述主要结果：中心 $\mathfrak{z}(\hat{\mathfrak{g}})$ 与 Langlands 对偶群 ${}^LG$ 所对应的圆盘上 oper 空间上的函数代数之间有一个典范同构。这个结果的证明构成本书的主要部分。我们将用它来描述临界水平下完备泛包络代数 $\widetilde{U}_{\kappa_c}(\hat{\mathfrak{g}})$ 的中心 $Z(\hat{\mathfrak{g}})$（定理 4.3.6）。最后，我们在命题 4.3.9 中证明：离开临界水平时，$\widetilde{U}_\kappa(\hat{\mathfrak{g}})$ 的中心是平凡的。

% [Tutor Note: 本章是全书的技术枢纽。p.101 的这条路线图应这样读：4.1 把“中心=射影联络”翻译成“中心=PGL_2-opers”，4.2 把 PGL_2 换成任意单李群 G，4.3 才给出真正的定理——中心等于 Langlands 对偶群 ^L G 的 opers 上的函数。注意对偶群在 4.2 中尚未出现，它只在 4.3 的定理陈述里登场；这正是本章最容易被读者忽略的转折。]

\begin{keypoints}
    \item \textbf{核心思想：}把 $\mathfrak{sl}_2$ 情形的“中心 $=$ 射影联络”提升为“中心 $=$ opers”，从而得到临界水平下顶点代数中心与完备泛包络代数中心的坐标无关描述。
    \item \textbf{数学工具：}平坦 $G$-丛与 $B$-约化、相对位置一形式、主分次与 $\mathfrak{sl}_2$-三元组、$N[[t]]$-规范等价类。
    \item \textbf{关键结论：}$\mathfrak{z}(\hat{\mathfrak{g}})\simeq\operatorname{Fun}\Op_{{}^LG}(D)$，$Z(\hat{\mathfrak{g}})\simeq\operatorname{Fun}\Op_{{}^LG}(D^\times)$，其中 ${}^LG$ 是 Langlands 对偶群。
    \item \textbf{记号提醒：}下标 ${}^LG$ 指对偶群 ${}^LG$ 的 opers，不是 $G$ 的 opers；$\mathfrak{z}$ 记顶点代数中心，$Z$ 记完备泛包络代数中心。
\end{keypoints}

\section{射影联络再探}

本节讨论射影联络的几种等价实现：射影结构与 $PGL_2$-opers。最后一种实现对我们最重要，因为它提示了如何把射影联络的概念（它在我们描述 $\mathfrak{g}=\mathfrak{sl}_2$ 的中心 $\mathfrak{z}(\hat{\mathfrak{g}})$ 时自然出现）推广到任意单李代数 $\mathfrak{g}$。

\begin{keypoints}
    \item \textbf{核心思想：}$X$ 上的射影联络、射影结构与 $PGL_2$-opers 是同一个对象的三种说法，最后一种才能推广。
    \item \textbf{数学工具：}二阶算子 $\Omega_X^{-1/2}\to\Omega_X^{3/2}$、Schwarz 导数、Möbius 图册、平坦 $PGL_2$-丛连带 $\mathbb{P}^1$-丛的整体截面。
    \item \textbf{关键结论：}oper 条件 $=$ 联络处处不保持 $B$-约化，即 (4.1-3) 中 $c(t)$ 处处非零；由此可化为标准形 (4.1-4)。
    \item \textbf{记号提醒：}$v_\alpha$ 是射影联络在坐标 $z_\alpha$ 下的系数，与 (4.1-4) 中的 $v(t)$ 是同一个量，按 (4.1-2) 变换。
\end{keypoints}

\subsection{射影结构}

§3.5.7 引入的射影联络概念不仅在圆盘上有意义，在 $\mathbb{C}$ 上任意光滑代数曲线 $X$ 上同样有意义。按定义，$X$ 上的射影联络是下列线丛的截面层之间的二阶微分算子：
\begin{equation}
    \rho:\Omega_X^{-1/2}\longrightarrow\Omega_X^{3/2},
    \tag{4.1-1}
\end{equation}
其主符号为 $1$，次主符号为 $0$。这里需要选取典范线丛 $\Omega_X$ 的平方根 $\Omega_X^{1/2}$。不过，正如 §3.5.7 所解释的，对应于 $\Omega_X^{1/2}$ 的不同选取，射影联络的空间之间有典范等同。\transnote{卡点：$\Omega_X^{1/2}$ 的选取看似使射影联络的定义依赖人为数据。背景：射影结构就是局部坐标到 $\mathbb{P}^1$ 的 Möbius 图册，它等价于主符号为 $1$、次主符号为 $0$ 的二阶算子 $\Omega_X^{-1/2}\to\Omega_X^{3/2}$，这是射影 Cartan 几何的平坦情形。结论：换用另一个平方根只把算子经线丛同构共轭，射影联络的空间与平方根选取无关，并且是以 $H^0(X,\Omega_X^2)$ 为模型的仿射空间。来源：E. Frenkel, Langlands Correspondence for Loop Groups, CUP 2007, p.101–102；nLab “projective structure”，https://ncatlab.org/nlab/show/projective+structure（访问 2026-11-23）。}

我们用 $\mathcal{P}roj(X)$ 记 $X$ 上射影联络的空间。

一个有用的观察是：$\mathcal{P}roj(X)$ 是以二次微分空间 $H^0(X,\Omega_X^2)$ 为模型的仿射空间。事实上，给定一个射影联络，即 (4.1-1) 中的二阶算子 $\rho$，再给定二次微分 $\omega\in H^0(X,\Omega_X^2)$，则 $\rho+\omega$ 仍是射影联络。此外，任给两个射影联络 $\rho,\rho'$，差 $\rho-\rho'$ 是零阶微分算子 $\Omega_X^{-1/2}\to\Omega_X^{3/2}$，等同于 $\Omega^2$ 的一个截面。可以证明，对任意光滑曲线（紧或非紧），空间 $\mathcal{P}roj(X)$ 非空（见 [Frenkel and Ben-Zvi 2004], §8.2.12）。因此它是 $H^0(X,\Omega_X^2)$-torsor。

局部地，在开解析子集 $U_\alpha\subset X$ 上取坐标 $z_\alpha$，并把线丛 $\Omega^{1/2}$ 平凡化。于是同 §3.5.7 一样，$\rho_\alpha=\rho|_{U_\alpha}$ 在此平凡化下可写成算子 $\partial_{z_\alpha}^2-v_\alpha(z_\alpha)$。同样如 §3.5.7，在重叠 $U_\alpha\cap U_\beta$ 上，设 $z_\alpha=\varphi_{\alpha\beta}(z_\beta)$，便有变换公式
\begin{equation}
    v_\beta(z_\beta)=v_\alpha(\varphi_{\alpha\beta}(z_\beta))\left(\frac{\partial\varphi_{\alpha\beta}}{\partial z_\beta}\right)^2-\frac{1}{2}\{\varphi_{\alpha\beta},z_\beta\}.
    \tag{4.1-2}
\end{equation}

% [Tutor Note: 核对 (4.1-2) 时最易出错的是三处约定：其一，$z_\alpha=\varphi_{\alpha\beta}(z_\beta)$ 的方向（若写成 $z_\beta=\varphi_{\beta\alpha}(z_\alpha)$，则求导因子要取倒数）；其二，Schwarz 导数前的 $-\tfrac12$ 而非 $-\tfrac16$ 或 $+\tfrac12$；其三，$\{\varphi,z\}=\varphi'''/\varphi'-\tfrac32(\varphi''/\varphi')^2$。三者中任一处符号或系数不符，$v_\alpha$ 就不能按 $\Omega^2$ 变换，射影联络的定义随之失效。]

下面我们把 $X$ 上的射影联络与另一类结构等同起来，这类结构现在给出定义。

$X$ 上的一个\textbf{射影图册}（projective chart）按定义是 $X$ 的一个开子集覆盖 $U_\alpha,\alpha\in A$，带局部坐标 $z_\alpha$，使得重叠 $U_\alpha\cap U_\beta$ 上的转移函数 $z_\beta=f_{\alpha\beta}(z_\alpha)$ 都是 Möbius 变换
\[
    f(z)=\frac{az+b}{cz+d},\qquad\begin{pmatrix}a&b\\c&d\end{pmatrix}\in PGL_2(\mathbb{C}).
\]
两个射影图册称为等价的，如果它们的并仍是射影图册。射影图册的等价类称为\textbf{射影结构}。

\begin{proposition}
    $X$ 上射影结构之集与 $X$ 上射影联络之集之间存在双射。
\end{proposition}

\begin{proof}
容易验证：对函数 $\varphi(z)$，Schwarz 导数 $\{\varphi,z\}$ 为 $0$ 当且仅当 $\varphi(z)$ 是 Möbius 变换。因此，给定射影结构，在每个图册 $U_\alpha$ 上指定二阶算子 $\partial_{z_\alpha}^2$，便定义了射影联络；由公式 (4.1-2)，这些算子相容地变换。

反之，给定射影联络，考虑微分方程
\[
    \left(\partial_{z_\alpha}^2-v_\alpha\right)\phi(z_\alpha)=0
\]
在 $U_\alpha$ 上的解空间。它二维，由 $\phi_{1,\alpha}$ 与 $\phi_{2,\alpha}$ 张成。取覆盖足够细，可以假设 $\phi_2$ 处处非零，且两解的朗斯基行列式（Wronskian）处处非零。于是定义
\[
    \mu_\alpha=\frac{\phi_{1,\alpha}}{\phi_{2,\alpha}}:U_\alpha\to\mathbb{C}.
\]
它良定义，且导数非零（因为朗斯基行列式非零）。故在原点附近它给出 $U_\alpha$ 上的复坐标。换一组基，这些 $\mu$ 之间显然由 Möbius 变换相联系。
\end{proof}

% [Tutor Note: “容易验证 Schwarz 导数消失当且仅当是 Möbius 变换”这一步，本书略去了必要的分析输入：二阶线性 ODE 的解空间二维，$\mu=\phi_1/\phi_2$ 的 Schwarz 导数为 $2v$（在坐标 $z_\alpha$ 下），故 $\{\mu,z_\alpha\}=2v_\alpha$；再由 (4.1-2) 的相容性可得 $\mu$ 在重叠处 Möbius 相关。可查 [Frenkel and Ben-Zvi 2004] §8.2 或 Ovsienko–Tabachnikov 的射影微分几何教材。]

\subsection{PGL$_2$-opers}

现在我们稍微改述射影结构的定义。这将给出思考射影联络的另一种方式，而它能推广到更一般的情形。

可以把 Möbius 变换看作群 $PGL_2(\mathbb{C})$ 的一个元素。于是，若 $X$ 上有射影结构，则对两个相交的图册 $U_\alpha$ 与 $U_\beta$，我们便对应了一个常映射
\[
    f_{\alpha\beta}:U_\alpha\cap U_\beta\longrightarrow PGL_2(\mathbb{C}).
\]
如 §1.2.4 所解释，这给出一个平坦 $PGL_2(\mathbb{C})$-丛的结构。等价地，如 §1.2.3 所述，它可表示成 $X$ 上一个全纯 $PGL_2$-丛 $\mathcal{F}$，连带一个全纯联络 $\nabla$（由于 $X$ 是曲线，$\nabla$ 自动平坦）。

然而我们还没有用到手头另有坐标 $z_\alpha$ 这一事实。群 $PGL_2(\mathbb{C})$ 自然作用在射影直线 $\mathbb{P}^1$ 上。作配 $\mathbb{P}^1$-丛
\[
    \mathbb{P}^1_{\mathcal{F}}=\mathcal{F}\times_{PGL_2(\mathbb{C})}\mathbb{P}^1
\]
于 $X$ 上。$\mathcal{F}$ 上的平坦结构在 $\mathbb{P}^1_{\mathcal{F}}$ 上诱导出平坦结构。于是我们有一个优选的系统，把 $\mathbb{P}^1_{\mathcal{F}}$ 的邻近纤维等同起来。现在，按 $\mathcal{F}$ 的定义，在每个 $U_\alpha$ 上丛变平凡，于是可用坐标 $z_\alpha$ 定义 $\mathbb{P}^1_{\mathcal{F}}|_{U_\alpha}=U_\alpha\times\mathbb{P}^1$ 的局部截面：让截面在点 $x$ 处取值 $z_\alpha(x)\in\mathbb{C}\subset\mathbb{P}^1$。

这其实是 $\mathbb{P}^1_{\mathcal{F}}$ 的整体截面，因为各坐标彼此之间的变换恰好就是我们用作丛在重叠开子集上的转移函数的那个 $PGL_2(\mathbb{C})$ 元素。由于 $z_\alpha$ 是坐标，它在每一点的导数非零。因此这些坐标给出 $\mathbb{P}^1$-丛 $\mathbb{P}^1_{\mathcal{F}}$ 的一个整体截面，且它（关于 $\mathbb{P}^1_{\mathcal{F}}$ 上的平坦联络）在每一点的导数都非零。

可以把这种结构想成：我们的黎曼面与射影直线之间的一组相容局部等同。“射影结构”之名即由此而来。

现在定义 $X$ 上的 \textbf{$PGL_2$-oper}：它是 $X$ 上一个平坦 $PGL_2(\mathbb{C})$-丛 $\mathcal{F}$，连带配 $\mathbb{P}^1$-丛 $\mathbb{P}^1_{\mathcal{F}}$ 的一个整体截面，使得该截面关于联络的导数处处非零。

于是我们看到，$X$ 上的射影结构（等价地，射影联络）给出 $X$ 上的 $PGL_2(\mathbb{C})$-oper。显然这个等同可逆：用截面定义局部坐标；平坦 $PGL_2$-丛的转移函数便定义了重叠处局部坐标的 Möbius 变换。因此，$X$ 上的射影结构（从而射影联络）与 $X$ 上的 $PGL_2$-opers 是同一回事。\transnote{卡点：为什么把射影结构改写成“平坦 $PGL_2$-丛 + $\mathbb{P}^1$-丛的整体截面”是有益的。背景：$\mathbb{P}^1=PGL_2/B$，故截面等价于 $B$-约化；截面关于平坦联络的导数非零，正是说该约化与联络横截，即相对位置落在开 $B$-轨道中。结论：射影结构、射影联络与 $PGL_2$-opers 三者典范等同，而“平坦 $G$-丛 + $B$-约化 + 横截”的形式可原样推广到任意单李群。来源：原书 p.103–104；A. Beilinson, V. Drinfeld, “Opers”, arXiv:math/0501398（2005）。}

现在的目标是把 $PGL_2$-oper 的概念推广到任意单李群。这将使我们能坐标无关地描述一般仿射 Kac--Moody 代数 $\hat{\mathfrak{g}}_{\kappa_c}$ 所对应的完备泛包络代数的中心。

要推广 oper 的概念，先要弄清 $\mathbb{P}^1$ 与群 $PGL_2(\mathbb{C})$ 的关系，以及截面导数的非零条件在说什么。

群 $PGL_2(\mathbb{C})$ 传递地作用在 $\mathbb{P}^1$ 上，而 $\mathbb{P}^1$ 可实现为 $\mathbb{C}^2=\operatorname{span}(e_1,e_2)$ 中的直线簇；直线 $\operatorname{span}(e_1)$ 的稳定子是 $PGL_2(\mathbb{C})$ 中上三角矩阵构成的 \textbf{Borel 子群} $B$。故 $\mathbb{P}^1$ 可表为齐性空间 $PGL_2(\mathbb{C})/B$。Borel 子群的概念容易推广到其他单李群，因此相应的齐性空间 $G/B$ 应是 $\mathbb{P}^1$ 的正确推广。

$\mathbb{P}^1$ 的元素可看作 $G$ 中 $B$ 的一个右陪集。于是，我们在 $\mathbb{P}^1_{\mathcal{F}}$ 的截面给出：在每个点 $x\in X$ 上方，纤维 $\mathcal{F}_x$ 中一个右 $B$-陪集。这意味着该截面给出主 $G$-丛的一个子丛，它是主 $B$-丛。

回忆主 $G$-丛 $\mathcal{F}$ 的 \textbf{$B$-约化}是 $B$-丛 $\mathcal{F}_B$ 连同同构 $\mathcal{F}\simeq\mathcal{F}_B\times_B G$。换言之，给定 $\mathcal{F}$ 的一个子丛，它是主 $B$-丛，且 $B$ 在其上的作用与从 $G$ 限制下来的作用一致。

\subsection{更多关于 oper 条件}

所以，一个 $PGL_2$-oper 给出平坦 $PGL_2$-丛 $(\mathcal{F},\nabla)$ 以及 $\mathcal{F}$ 到 Borel 子群 $B$ 的一个约化 $\mathcal{F}_B$。数据 $\mathcal{F}_B$ 等价于 $\mathbb{P}^1_{\mathcal{F}}$ 的一个截面。现在须把 $PGL_2$-opers 定义中出现的、关于 $\mathbb{P}^1_{\mathcal{F}}$ 截面导数的非零条件，用 $\mathcal{F}_B$ 表出。这个条件基本上是说：$\mathcal{F}$ 上的平坦联络处处不保持 $B$-子丛 $\mathcal{F}_B$。下面更精确地说明。

如 §1.2.4 所述，若在开子集 $U\subset X$ 上取局部坐标 $t$，并把全纯 $PGL_2$-丛 $\mathcal{F}$ 平凡化，则联络 $\nabla$ 给出一个一阶微分算子
\begin{equation}
    \nabla_{\partial_t}=\partial_t+\begin{pmatrix}a(t)&b(t)\\c(t)&d(t)\end{pmatrix},
    \tag{4.1-3}
\end{equation}
其中矩阵在李代数 $\mathfrak{g}=\mathfrak{sl}_2(\mathbb{C})$ 中，即 $a(t)+d(t)=0$。此算子不典范，改变 $\mathcal{F}$ 的平凡化时它按规范变换变化（它也随坐标 $t$ 的改变而变换）。

$PGL_2$-oper 上的条件可以改述为：取 $U$ 上 $\mathcal{F}$ 的一个关于该联络为水平的局部平凡化，使 $\nabla_{\partial_t}=\partial_t$。则 $\mathbb{P}^1_{\mathcal{F}}$ 的截面可表示为函数 $f:U\to\mathbb{P}^1$，而条件是它对 $t$ 的导数处处非零。

% [Tutor Note: (4.1-3) 把 PGL_2-丛的联络写成 sl_2 中的矩阵，这里有一个易被忽略的细节：PGL_2 的丛不一定是 SL_2 的丛，故矩阵整体只定义到 ±I；(4.1-4) 之后用对角矩阵 (c^{1/2},c^{-1/2}) 作规范变换时，正因如此才说“在 PGL_2 中有定义”。核对符号时请连带核对这一处约定。]

上述描述正确，但不太有用，因为一般很难显式找到水平平凡化：那需要解一个微分方程组，可能不容易。

更实用的做法是改用由 $\mathcal{F}_B$ 在 $U$ 上某个平凡化诱导的 $\mathcal{F}$ 的平凡化。此时 $f$ 成为常映射，其值为 $B$ 在 $\mathbb{P}^1=PGL_2/B$ 中的陪集。但联络现在取 (4.1-3) 的形式，矩阵元是非平凡函数。这些函数只在规范变换意义下确定。不过，由于 $\mathcal{F}$ 的平凡化由 $\mathcal{F}_B$ 的平凡化诱导，这种不确定性只由取值为 $B$ 的规范变换构成。我们想对 $\nabla_{\partial_t}$ 加的条件，都应在这些 $B$-值规范变换（以及坐标 $t$ 的改变）下不变。例如，保持我们的 $B$-约化的联络恰是 $c(t)=0$ 的那些；这显然是 $B$-值规范变换下不变的条件。

相反，我们的条件是：由 $\mathbb{P}^1_{\mathcal{F}}$ 的截面借助 $\mathcal{F}_B$ 的平凡化所得的函数 $f:U\to\mathbb{P}^1$，它关于 $\nabla_{\partial_t}$ 的导数处处非零。注意此导数在点 $x\in U$ 的值其实是 $\mathbb{P}^1$ 在 $f(x)$ 处的切向量。由于 $f(x)=B\in PGL_2/B=\mathbb{P}^1$，$f(x)$ 处的切空间自然等同于 $\mathfrak{sl}_2/\mathfrak{b}$。因此该导数由公式 (4.1-3) 中的矩阵元 $c(t)$ 表示。于是条件就是 $c(t)$ 处处非零。

满足此性质的联络可用 $B$-值规范变换化为标准形。首先用对角矩阵
\[
    \begin{pmatrix}c(t)^{1/2}&0\\0&c(t)^{-1/2}\end{pmatrix}
\]
作规范变换（注意它在 $PGL_2$ 中有定义），把联络算子化为 $c(t)=1$ 的 (4.1-3) 形式。接着用上三角变换
\[
    \begin{pmatrix}1&-a\\0&1\end{pmatrix}\left(\partial_t+\begin{pmatrix}a&b\\1&-a\end{pmatrix}\right)\begin{pmatrix}1&a(t)\\0&1\end{pmatrix}
    =\partial_t+\begin{pmatrix}0&b+a^2+\partial_t a\\1&0\end{pmatrix}.
\]
结果得到形如
\begin{equation}
    \nabla=\partial_t+\begin{pmatrix}0&v(t)\\1&0\end{pmatrix}
    \tag{4.1-4}
\end{equation}
的联络算子。

由上述构造可见，存在唯一的 $B$-值函数（上述对角矩阵与上三角矩阵之积），把给定的、$c(t)$ 处处非零的算子 (4.1-3) 化为这种形式。于是我们用尽了 $B$-值规范变换的全部自由来把联络化为该形式。

这意味着：在小的开子集 $U\subset X$ 上（关于坐标 $t$），或在点 $x\in X$ 周围的圆盘 $D_x$ 上（关于某个形式坐标 $t$ 的选取），$PGL_2$-opers 的空间等同于形如 (4.1-4) 的算子空间，其中 $v(t)$ 是 $U$ 上的函数（相应地，$v(t)\in\mathbb{C}[[t]]$）。

由矩阵微分方程理论可知，微分方程 $\nabla\Phi(t)=0$ 等价于二阶微分方程 $(\partial_t^2-v(t))\Phi(t)=0$。因此可把 opers 空间等同于算子 $\partial_t^2-v(t)$ 的空间。它们看起来像射影联络，但要把它们与射影联络等同，还须验证：在坐标变换下它们确实按微分算子 $\Omega_X^{-1/2}\to\Omega_X^{3/2}$ 变换。

% [Tutor Note: 上三角变换中第三个矩阵写作 $\begin{pmatrix}1&a(t)\\0&1\end{pmatrix}$ 而第二个矩阵里用的是 $a$，原文在记号上不统一（应为同一个 $a(t)$）。这不影响结果，但抄录公式时容易误以为 $a$ 与 $a(t)$ 是两个不同的函数。]

为此，把平坦 $PGL_2$-丛 $\mathcal{F}$ 实现为秩二平坦向量丛模去与平坦线丛作张量积后的等价类。利用与平坦线丛作张量积的自由，可在该类中选取代表，使其行列式为平凡线丛带平凡联络。该代表也记作 $\mathcal{F}$。它不唯一，但可与平方为平凡的任何平坦线丛 $(\mathcal{L},\nabla)$ 作张量积。$B$-约化 $\mathcal{F}_B$ 给出线子丛 $\mathcal{F}_1\subset\mathcal{F}$，它只在上面的平坦线丛 $(\mathcal{L},\nabla)$ 的张量积意义下确定。

联络 $\nabla$ 是映射 $\mathcal{F}\otimes\mathcal{T}_X\to\mathcal{F}$，或等价地 $\mathcal{F}\to\mathcal{F}\otimes\Omega_X$，因为 $\Omega_X$ 与 $\mathcal{T}_X$ 是对偶线丛。oper 条件在这个观点下是什么意思？它恰是说：复合
\[
    \mathcal{F}_1\xrightarrow{\ \nabla\ }\mathcal{F}\otimes\Omega_X\longrightarrow(\mathcal{F}/\mathcal{F}_1)\otimes\Omega_X
\]
是线丛的同构（它对应 (4.1-3) 中出现的函数 $c(t)$）。

由于已把 $\det\mathcal{F}\simeq\mathcal{F}_1\otimes(\mathcal{F}/\mathcal{F}_1)$ 与 $\mathcal{O}_X$ 等同，这意味着 $(\mathcal{F}_1)^{\otimes2}\simeq\Omega_X$，从而 $\mathcal{F}_1\simeq\Omega_X^{1/2}$。这也推出 $\mathcal{F}/\mathcal{F}_1\simeq\Omega_X^{-1/2}$。这两个同构都不依赖 $\Omega_X^{1/2}$ 的选取，因为 $\mathcal{F}_1$ 本来就只在模去与平凡线丛的平方根作张量积的意义下有定义。取定一个平方根 $\Omega_X^{1/2}$，便可写
\[
    0\longrightarrow\Omega_X^{1/2}\longrightarrow\mathcal{F}\longrightarrow\Omega_X^{-1/2}\longrightarrow 0.
\]

暂设 $X$ 是亏格 $g$ 的射影曲线。则
\[
    \operatorname{Ext}^1(\Omega_X^{-1/2},\Omega_X^{1/2})\cong H^1(X,\Omega_X)\cong H^0(X,\mathcal{O}_X)^*\cong\mathbb{C},
\]
因此 $\Omega_X^{-1/2}$ 被 $\Omega_X^{1/2}$ 的扩张至多有两个同构类：分裂的与非分裂的。若扩张分裂，则线丛 $\Omega_X^{1/2}$ 上有诱导联络；但该丛次数为 $g-1$，当 $g\ne1$ 时不可能带平坦联络。故 $g\ne1$ 时，对应 $H^1(X,\Omega_X)$ 中非零类的 oper 丛唯一。（$g=1$ 时两种扩张彼此同构，故 oper 丛仍唯一。）

% [Tutor Note: 这一段用到 Serre 对偶 $H^1(X,\Omega_X)\cong H^0(X,\mathcal{O}_X)^*$（对曲线成立），以及“平坦线丛次数为 0”这一事实。注意 $g=1$ 处作者只给了一句括号说明：此时 $\deg\Omega_X^{1/2}=0$，分裂扩张给出的联络存在，但两类扩张本身同构，故 oper 丛仍唯一。核对时请勿把“唯一 oper 丛”误读成“唯一 oper”。]

对任意光滑曲线 $X$，由联络的形式可得：若 $\mathcal{F}$ 有水平截面，则其到 $\Omega_X^{-1/2}$ 的投影是某个二阶算子的解，该算子就是我们的 $\partial_t^2-v(t)$。同样清楚的是，此算子作用 $\Omega_X^{-1/2}\to\Omega_X^{3/2}$，故确是射影联络。由于平坦丛与射影联络都由其水平截面（即解）唯一确定，我们得到：$PGL_2$-opers 与射影联络一一对应。

我们兜了一个整圈。从射影联络出发，改述为射影结构，再译成 $PGL_2$-opers 的说法，最后利用 oper 联络的标准代表元发现它自然地对应射影联络。于是我们找到同一个对象的三种化身：射影联络、射影结构与 $PGL_2$-opers。我们将把最后这个概念推广到以任意单李代数替代 $\mathfrak{sl}_2$。

\section{一般单李代数的 opers}

本节引入与任意单李代数（等价地，伴随型单李群）相联系的 opers。下一节将看到：opers 之于顶点代数 $V_{\kappa_c}(\mathfrak{g})$ 的中心，正如射影联络之于 $V_{\kappa_c}(\mathfrak{sl}_2)$ 的中心。但故事有一个重要的转折：中心同构于圆盘上 oper 空间上的函数代数，而这里的 oper 不是联系于 $G$，而是联系于 $G$ 的 Langlands 对偶群 ${}^LG$。对偶群的出现，从局部 Langlands 对应的观点看非常重要。

\begin{keypoints}
    \item \textbf{核心思想：}用任意单李群 $G$ 的 Borel 约化代替 $\mathbb{P}^1$，oper 条件即要求相对位置一形式落在开 $B$-轨道 $\mathbf{O}\subset(\mathfrak{g}/\mathfrak{b})$ 中。
    \item \textbf{数学工具：}相对位置一形式 $\nabla/\mathcal{F}_B$、主分次、指数 $d_i$、$\mathfrak{sl}_2$-三元组、$N[[t]]$-规范等价类。
    \item \textbf{关键结论：}每个 oper 唯一地写成 $\partial_t+p_{-1}+\sum_{j}v_j(t)p_j$；$\Op_G(D_x)\simeq\mathcal{P}roj(D_x)\times V^{\mathrm{abs}}_{>1,x}$。
    \item \textbf{记号提醒：}$V^{\mathrm{can}}$ 依赖 $\mathfrak{sl}_2$-三元组的选取，$V^{\mathrm{abs}}$ 才与选取无关；$p_{-1}=\sum f_i$。
\end{keypoints}

\subsection{opers 的定义}

设 $G$ 是伴随型单代数群，$B$ 是其 Borel 子群，$N=[B,B]$ 是其幂幺根，对应的李代数为 $\mathfrak{n}\subset\mathfrak{b}\subset\mathfrak{g}$。

于是 $\mathfrak{g}$ 是单李代数，因而有 Cartan 分解
\[
    \mathfrak{g}=\mathfrak{n}_-\oplus\mathfrak{h}\oplus\mathfrak{n}_+.
\]
我们选取 $\mathfrak{n}_+$（相应地 $\mathfrak{n}_-$）的生成元 $e_1,\dots,e_\ell$（相应地 $f_1,\dots,f_\ell$）。有 $\mathfrak{n}_{\alpha_i}=\mathbb{C}e_i$，$\mathfrak{n}_{-\alpha_i}=\mathbb{C}f_i$，见附录 A.3。取 $\mathfrak{b}=\mathfrak{h}\oplus\mathfrak{n}_+$ 为 $B$ 的李代数，则 $\mathfrak{n}$ 是 $N$ 的李代数。下文用记号 $\mathfrak{n}$ 表示 $\mathfrak{n}_+$。

对我们有用的是这样约定：只选取 $\mathfrak{n}=\mathfrak{n}_+$ 与 $\mathfrak{b}=\mathfrak{h}\oplus\mathfrak{n}_+$，而不选取 $\mathfrak{h}$ 或 $\mathfrak{n}_-$。于是我们有一个抽象的 Cartan 李代数 $\mathfrak{h}=\mathfrak{b}/\mathfrak{n}$，但没有 $\mathfrak{h}$ 到 $\mathfrak{b}$ 的嵌入。有时我们会选取这样的嵌入，此时每个选取也给出下幂零子代数 $\mathfrak{n}_-$（定义为此嵌入所对应的负根向量张成的空间）。同样，我们有子群 $N\subset B\subset G$ 与抽象 Cartan 群 $H=G/B$，但没有分裂 $H\hookrightarrow B$。凡用到分裂 $\mathfrak{h}\hookrightarrow\mathfrak{b}$（等价地 $H\hookrightarrow B$）之处，我们都会说明这个选取如何影响讨论。\transnote{卡点：这里原书把抽象 Cartan 群印作 $H=G/B$，但 $G/B$ 是旗簇，并不是群。背景：抽象 Cartan 群应为 Borel 群模去其幂幺根，即 $H=B/N$；同一页下文正用它来断言 $\mathbf{O}$ 是 $H$-torsor。结论：译文保留原书这一处笔误，读者应按 $H=B/N$ 理解，后文论证不受影响。来源：依据原书 p.108（PDF p.122）同页上下文。}

设 $[\mathfrak{n},\mathfrak{n}]^\perp\subset\mathfrak{g}$ 是 $[\mathfrak{n},\mathfrak{n}]$ 关于非退化不变内积 $\kappa_0$ 的正交补。有
\[
    [\mathfrak{n},\mathfrak{n}]^\perp/\mathfrak{b}\simeq\bigoplus_{i=1}^{\ell}\mathfrak{n}_{-\alpha_i}.
\]
显然群 $B$ 作用在 $\mathfrak{n}^\perp/\mathfrak{b}$ 上。第一个观察是：存在开的 $B$-轨道 $\mathbf{O}\subset\mathfrak{n}^\perp/\mathfrak{b}\subset\mathfrak{g}/\mathfrak{b}$，它由在每个子空间 $\mathfrak{n}_{-\alpha_i}$ 上的投影都非零的向量组成。该轨道也可描述为：把任意可能的子代数 $\mathfrak{n}_-$ 的生成元 $f_i,\,i=1,\dots,\ell$ 投影到 $\mathfrak{g}/\mathfrak{b}$ 后求和，所得向量的 $B$-轨道。$B$ 在 $\mathbf{O}$ 上的作用经由 $H=B/N$ 的作用分解。后者单传递地作用，使 $\mathbf{O}$ 成为 $H$-torsor（见 §3.5.3）。

% [Tutor Note: 原文此处先说 $[\mathfrak{n},\mathfrak{n}]^\perp\subset\mathfrak{g}$ 及其商 $[\mathfrak{n},\mathfrak{n}]^\perp/\mathfrak{b}$，紧接着却写“$B$ 作用在 $\mathfrak{n}^\perp/\mathfrak{b}$ 上”、“$\mathbf{O}\subset\mathfrak{n}^\perp/\mathfrak{b}$”。按不变内积的性质，$\mathfrak{n}^\perp=\mathfrak{b}$，故 $\mathfrak{n}^\perp/\mathfrak{b}=0$，这里的 $\mathfrak{n}^\perp$ 应是 $[\mathfrak{n},\mathfrak{n}]^\perp$ 的简写（或笔误）。全节凡是出现 $\mathfrak{n}^\perp/\mathfrak{b}$ 处都按 $[\mathfrak{n},\mathfrak{n}]^\perp/\mathfrak{b}$ 理解，否则开轨道将是空集。]

设 $X$ 是光滑曲线，$x\in X$。仍用 $\mathcal{O}_x$ 记完备局部环，$\mathcal{K}_x$ 记其分式域。环 $\mathcal{O}_x$ 同构于 $\mathbb{C}[[t]]$，但不典范。于是 $D_x=\operatorname{Spec}\mathcal{O}_x$ 是不带坐标的圆盘，$D_x^\times=\operatorname{Spec}\mathcal{K}_x$ 是相应的去心圆盘。

现设给定光滑曲线 $X$（或 $D_x$，或 $D_x^\times$）上的主 $G$-丛 $\mathcal{F}$，连带联络 $\nabla$（自动平坦）以及到 Borel 子群 $B$ 的约化 $\mathcal{F}_B$。定义 $\nabla$ 与 $\mathcal{F}_B$ 的\textbf{相对位置}（即 $\nabla$ 保持 $\mathcal{F}_B$ 的失败程度）如下。局部取 $\mathcal{F}$ 上任意保持 $\mathcal{F}_B$ 的平坦联络 $\nabla'$，考虑差 $\nabla-\nabla'$，它是 $\mathfrak{g}_{\mathcal{F}_B}\otimes\Omega_X$ 的截面。把它投影到 $(\mathfrak{g}/\mathfrak{b})_{\mathcal{F}_B}\otimes\Omega_X$。显然，所得的 $(\mathfrak{g}/\mathfrak{b})_{\mathcal{F}_B}\otimes\Omega_X$ 的局部截面不依赖 $\nabla'$ 的选取。这些截面拼接起来，定义 $X$ 上一个整体的 $(\mathfrak{g}/\mathfrak{b})_{\mathcal{F}_B}$-值一形式，记作 $\nabla/\mathcal{F}_B$。

设 $X$ 是光滑曲线，或 $D_x$，或 $D_x^\times$。设给定 $X$ 上的主 $G$-丛 $\mathcal{F}$、$\mathcal{F}$ 上的联络 $\nabla$ 与 $B$-约化 $\mathcal{F}_B$。若一形式 $\nabla/\mathcal{F}_B$ 取值于 $\mathbf{O}_{\mathcal{F}_B}\subset(\mathfrak{g}/\mathfrak{b})_{\mathcal{F}_B}$，则称 $\mathcal{F}_B$ 与 $\nabla$ \textbf{横截}。注意 $\mathbf{O}$ 是 $\mathbb{C}^\times$-不变的，故 $\mathbf{O}\otimes\Omega_X$ 是 $(\mathfrak{g}/\mathfrak{b})_{\mathcal{F}_B}\otimes\Omega_X$ 的良定义子集。

现在，$X$ 上的 \textbf{$G$-oper} 按定义是三元组 $(\mathcal{F},\nabla,\mathcal{F}_B)$，其中 $\mathcal{F}$ 是 $X$ 上的主 $G$-丛，$\nabla$ 是 $\mathcal{F}$ 上的联络，$\mathcal{F}_B$ 是 $\mathcal{F}$ 的 $B$-约化，且 $\mathcal{F}_B$ 与 $\nabla$ 横截。

这个定义属于 A. Beilinson 与 V. Drinfeld [1997]（当 $X$ 是去心圆盘时，opers 更早由 V. Drinfeld 与 V. Sokolov 在 [1985] 中引入）。

显然，当 $G=PGL_2$ 时，我们回到 §4.1.2 中 $PGL_2$-oper 的定义。\transnote{卡点：oper 定义中的横截条件为何写成相对位置 $\nabla/\mathcal{F}_B$ 取值于开轨道 $\mathbf{O}$；原书此处又把 $[\mathfrak{n},\mathfrak{n}]^\perp/\mathfrak{b}$ 简写成 $\mathfrak{n}^\perp/\mathfrak{b}$（因 $\mathfrak{n}^\perp=\mathfrak{b}$，按字面读 $\mathbf{O}$ 为空集）。背景：$[\mathfrak{n},\mathfrak{n}]^\perp/\mathfrak{b}\simeq\bigoplus_i\mathfrak{n}_{-\alpha_i}$，横截即 $\nabla/\mathcal{F}_B$ 在每个 $\mathfrak{n}_{-\alpha_i}$ 方向的分量都非零。结论：按 $[\mathfrak{n},\mathfrak{n}]^\perp/\mathfrak{b}$ 理解，当 $G=PGL_2$ 时正回到 §4.1.3 的 $c(t)$ 处处非零。来源：原书 p.107–108；A. Beilinson, V. Drinfeld, “Opers”, arXiv:math/0501398（2005）。}

\subsection{实现为规范等价类}

等价地，横截条件可改述为：若取 $\mathcal{F}_B$ 的局部平凡化与局部坐标 $t$，则联络具有形式
\begin{equation}
    \nabla=\partial_t+\sum_{i=1}^{\ell}\psi_i(t)f_i+\mathbf{v}(t),
    \tag{4.2-1}
\end{equation}
其中每个 $\psi_i(t)$ 是处处非零函数，$\mathbf{v}(t)$ 是 $\mathfrak{b}$-值函数。其证明与 §4.1.3 中 $G=PGL_2$ 的情形完全相同。

若改变 $\mathcal{F}_B$ 的平凡化，则该算子按相应的 $B$-值规范变换变换（见 §1.2.4）。这一观察使我们能更具体地描述圆盘 $D_x=\operatorname{Spec}\mathcal{O}_x$ 与去心圆盘 $D_x^\times=\operatorname{Spec}\mathcal{K}_x$ 上的 opers。同样的推理对任意曲线上任意足够小的解析子集 $U$（带局部坐标 $t$），或带 étale 坐标的 Zariski 开子集也成立。为确定起见，以下考虑底为 $D_x$ 的情形。

在 $D_x$ 上取坐标 $t$，即同构 $\mathcal{O}_x\simeq\mathbb{C}[[t]]$。于是把 $D_x$ 与 $D=\operatorname{Spec}\mathbb{C}[[t]]$ 等同。$D$ 上 $G$-opers 的空间 $\Op_G(D)$ 是形如 (4.2-1) 的所有算子（其中 $\psi_i(t)\in\mathbb{C}[[t]]$，$\psi_i(0)\ne0$，$i=1,\dots,\ell$，$\mathbf{v}(t)\in\mathfrak{b}[[t]]$）构成的空间，模去规范变换群 $B[[t]]$ 的作用：
\[
    g\cdot(\partial_t+A(t))=\partial_t+gA(t)g^{-1}-g^{-1}\partial_t g.
\]

取同态 $B\to H$ 的分裂 $\iota:H\to B$。则 $B$ 成为半直积 $B=H\ltimes N$。$B$-轨道 $\mathbf{O}$ 是 $H$-torsor，故可用 $H$-值规范变换把所有函数 $\psi_i(t)$ 变为 $1$。换言之，存在唯一的 $H[[t]]$ 元素，即 $\prod_{i=1}^{\ell}\check{\omega}_i(\psi_i(t))$（其中 $\check{\omega}_i:\mathbb{C}^\times\to H$ 是 $G$ 的第 $i$ 个基本余权），使得相应的规范变换把联络算子化为
\begin{equation}
    \nabla=\partial_t+\sum_{i=1}^{\ell}f_i+\mathbf{v}(t),\qquad\mathbf{v}(t)\in\mathfrak{b}[[t]].
    \tag{4.2-2}
\end{equation}
剩下的是 $N$-值规范变换群。于是得到：$\Op_G(D)$ 等于形如 (4.2-2) 的算子空间 $\widetilde{\Op}_G(D)$ 模去群 $N[[t]]$ 的规范变换作用：
\[
    \Op_G(D)=\widetilde{\Op}_G(D)/N[[t]].
\]

这给出圆盘上 opers 空间作为规范等价类的一个非常具体的实现。\transnote{卡点：为何每个 oper 都可由 (4.2-2) 形算子唯一代表，$H[[t]]$ 与 $N[[t]]$ 各消去多少自由度。背景：取定分裂 $B=H\ltimes N$ 后，$H$ 单传递地作用于开轨道 $\mathbf{O}$，故存在唯一的 $H[[t]]$ 元素 $\prod_i\check\omega_i(\psi_i(t))$ 把所有 $\psi_i(t)$ 化为 $1$；剩下的 $N[[t]]$ 作用由引理 4.2.2 保证自由。结论：$\Op_G(D)=\widetilde{\Op}_G(D)/N[[t]]$，而 $\widetilde{\Op}_G(D)$ 中的元素唯一地写成 $\partial_t+p_{-1}+\mathbf{c}(t)$（$\mathbf{c}(t)\in V[[t]]$）。来源：原书 p.110、113–114 与 [Drinfeld and Sokolov 1985]（本书文献表）。}

% [Tutor Note: 此处“用 H-值规范变换把 $\psi_i(t)$ 全变成 1”之所以唯一，是因为 $\mathbf{O}$ 是 $H$-torsor（$H$ 单传递地作用其上），而不是仅仅有传递性。$\prod_i\check\omega_i(\psi_i(t))$ 就是把向量 $(\psi_1,\dots,\psi_\ell)$ 所在的 $H$-轨道送回基点的那唯一元素。核对 (4.2-2) 时可顺带验证 $[\check\rho,e_i]=e_i$、$[\check\rho,f_i]=-f_i$：正是它保证 $\psi_i\equiv1$ 的形式在 $H$-规范变换下不变。]

\subsection{坐标变换的作用}

上述公式中我们用了圆盘 $D_x$ 上的一个特定坐标 $t$（并用它把 $D_x$ 与 $D$ 等同）。由于目标是用坐标无关的方式陈述所有结果，需要弄清：若换用另一个坐标，上面引入的规范等价类如何改变。

设 $s$ 是 $D_x$ 上的另一坐标，$t=\varphi(s)$。用新坐标，算子 (4.2-2) 变为
\[
    \nabla_{\partial_t}=\nabla_{\varphi'(s)^{-1}\partial_s}=\varphi'(s)^{-1}\partial_s+\sum_{i=1}^{\ell}f_i+\mathbf{v}(\varphi(s)).
\]
于是
\[
    \nabla_{\partial_s}=\partial_s+\varphi'(s)\sum_{i=1}^{\ell}f_i+\varphi'(s)\cdot\mathbf{v}(\varphi(s)).
\]
要把它化回 (4.2-2) 的形式，需作由 $\check{\rho}(\varphi'(s))$ 给出的规范变换，其中 $\check{\rho}:\mathbb{C}^\times\to H$ 是 $H$ 的单参数子群，等于 $G$ 的基本余权之和，$\check{\rho}=\sum_{i=1}^{\ell}\check{\omega}_i$。（这里再次取同态 $B\to H$ 的分裂 $\iota:H\to B$。）把 $\check{\rho}$ 看作李代数 $\mathfrak{h}=\operatorname{Lie}H$ 的元素，有 $[\check{\rho},e_i]=e_i$，$[\check{\rho},f_i]=-f_i$（见附录 A.3）。因此
\begin{align}
    &\check{\rho}(\varphi'(s))\cdot\left(\partial_s+\varphi'(s)\sum_{i=1}^{\ell}f_i+\varphi'(s)\cdot\mathbf{v}(\varphi(s))\right)\nonumber\\
    &\qquad=\partial_s+\sum_{i=1}^{\ell}f_i+\varphi'(s)\check{\rho}(\varphi'(s))\cdot\mathbf{v}(\varphi(s))\cdot\check{\rho}(\varphi'(s))^{-1}-\check{\rho}\left(\frac{\varphi''(s)}{\varphi'(s)}\right).
    \tag{4.2-3}
\end{align}

上述公式定义了群 $\operatorname{Aut}\mathcal{O}$ 在标准圆盘 $D=\operatorname{Spec}\mathbb{C}[[t]]$ 上 opers 空间 $\Op_G(D)$ 上的作用。\transnote{卡点：(4.2-3) 中的规范变换为何恰是 $\check\rho(\varphi'(s))$，末项 $\check\rho(\varphi''(s)/\varphi'(s))$ 从何而来。背景：$\check\rho=\sum_i\check\omega_i$ 是主分次的生成元，$[\check\rho,e_i]=e_i$、$[\check\rho,f_i]=-f_i$，故 $\check\rho(\varphi')$ 的共轭把 $\varphi'\sum_i f_i$ 的系数化为 $1$。结论：末项来自 $\check\rho(\varphi')^{-1}\partial_s\check\rho(\varphi')$，它正是 (4.1-2) 中 Schwarz 导数项在一般 $\mathfrak{g}$ 中的对应物。来源：原书 p.110–111 与附录 A.3。}

对于光滑曲线 $X$ 上点 $x$ 周围的圆盘 $D_x$，现在可把 $\Op_G(D_x)$ 定义为 $\Op_G(D)$ 沿 $\operatorname{Aut}\mathcal{O}$-torsor $\mathcal{A}ut_x$ 的扭（见 §3.5.2）。

特别地，上述公式使我们能确定 $H$-丛 $\mathcal{F}_H=\mathcal{F}_B\times_B H=\mathcal{F}_B/N$ 的结构。先描述 $H$-丛的一般构造，它在任意光滑曲线、圆盘或去心圆盘上都适用。

设 $\check{P}$ 是余特征标 $\mathbb{C}^\times\to H$ 的格。由于 $G$ 是相应于 $G$ 的伴随型群，该格自然等同于 Cartan 代数 $\mathfrak{h}$ 的整余权格，由 $\check{\omega}_i,\,i=1,\dots,\ell$ 张成。设 $P$ 是特征标 $H\to\mathbb{C}^\times$ 的格。我们有自然配对 $\langle\cdot\,,\cdot\rangle:P\times\check{P}\to\mathbb{Z}$，由特征标与余特征标复合得到，给出同态 $\mathbb{C}^\times\to\mathbb{C}^\times$，对应一个整数。

给定 $X$ 上的 $H$-丛 $\mathcal{F}_H$，对每个 $\lambda\in P$ 有配线丛 $\mathcal{L}_\lambda=\mathcal{F}_H\times_H\mathbb{C}_\lambda$，其中 $\mathbb{C}_\lambda$ 是 $H$ 对应于 $\lambda$ 的一维表示。数据 $\mathcal{F}_H$ 等价于线丛 $\{\mathcal{L}_\lambda,\lambda\in P\}$ 连同满足显然结合性条件的同构 $\mathcal{L}_\lambda\otimes\mathcal{L}_\mu\simeq\mathcal{L}_{\lambda+\mu}$。

% [Tutor Note: 原文“由于 $G$ 是相应于 $G$ 的伴随型群”显然是笔误，应为“相应于 $\mathfrak{g}$ 的伴随型群”：$G$ 是以 $\mathfrak{g}$ 为李代数的伴随型群，故其（余特征标）格恰是整余权格（若 $G$ 单连通则不成立）。这条辨识在 §4.3 中描述 ${}^LG$-opers 时还要再用。]

现在给定 $\check{\mu}\in\check{P}$，令 $\mathcal{L}_\lambda=\Omega_X^{\langle\lambda,\check{\mu}\rangle}$ 来定义这样的数据。把对应的 $X$ 上 $H$-丛记作 $\Omega^{\check{\mu}}$。它不过是 $\mathbb{C}^\times$-丛 $\Omega^\times$ 沿余特征标 $\check{\mu}$ 给出的同态 $\mathbb{C}^\times\to H$ 的前推。

更具体地，$H$-丛 $\Omega^{\check{\mu}}$ 可描述如下：对 $X$ 的开子集上局部坐标 $t$ 的每个选取，有 $\Omega_X$ 由截面 $dt$ 生成的平凡化，从而每个线丛 $\mathcal{L}_\lambda$ 也有平凡化。若换用另一坐标 $s$，$t=\varphi(s)$，则两个平凡化相差转移函数 $\langle\lambda,\check{\mu}(\varphi'(s))\rangle$。换言之，$\Omega^{\check{\mu}}$ 的转移函数是 $\check{\mu}(\varphi'(s))$。这些转移函数给出定义 $H$-丛 $\Omega^{\check{\mu}}$ 的另一种方式。

\begin{lemma}
    $H$-丛 $\mathcal{F}_H=\mathcal{F}_B\times_B H=\mathcal{F}_B/N$ 同构于 $\Omega^{\check{\rho}}$。
\end{lemma}

\begin{proof}
由公式 (4.2-3)（坐标变换对 opers 的作用）可得：若把 $D_x$ 上的坐标 $t$ 换成坐标 $s$（$t=\varphi(s)$），则得到 $H$-丛 $\mathcal{F}_H$ 的一个新平凡化，它通过转移函数 $\check{\rho}(\varphi'(s))$ 与旧平凡化相联系。这恰意味着 $\mathcal{F}_H\simeq\Omega^{\check{\rho}}$。
\end{proof}

\subsection{标准代表元}

本节在 (4.2-2) 形联络的 $N[[t]]$-规范类中找出标准代表元。

为此注意，算子 $\operatorname{ad}\check{\rho}$ 在 $\mathfrak{g}$ 上定义一个分次，称为\textbf{主分次}；关于它有直和分解 $\mathfrak{g}=\bigoplus_i\mathfrak{g}_i$。特别地，$\mathfrak{b}=\bigoplus_{i\ge0}\mathfrak{b}_i$，其中 $\mathfrak{b}_0=\mathfrak{h}$。

令
\[
    p_{-1}=\sum_{i=1}^{\ell}f_i.
\]
算子 $\operatorname{ad}p_{-1}$ 对所有 $i\ge0$ 从 $\mathfrak{b}_{i+1}$ 到 $\mathfrak{b}_i$ 单射。因此对每个 $i\ge0$ 可找到子空间 $V_i\subset\mathfrak{b}_i$，使得 $\mathfrak{b}_i=[p_{-1},\mathfrak{b}_{i+1}]\oplus V_i$。众所周知，$V_i\ne0$ 当且仅当 $i$ 是 $\mathfrak{g}$ 的\textbf{指数}，此时 $\dim V_i$ 等于该指数的重数。特别地，$V_0=0$。

令 $V=\bigoplus_{i\in E}V_i\subset\mathfrak{n}$，其中 $E=\{d_1,\dots,d_\ell\}$ 是计重数的 $\mathfrak{g}$ 的指数集合。这些指数正是我们在 §2.1.1 中遇到过的那一组：它们等于 $U(\mathfrak{g})$ 的中心的生成元的次数减 $1$。

注意除 $\mathfrak{g}=D_{2n}$、$d_n=2n$（此时重数为 $2$）的情形外，每个指数的重数都是 $1$。\transnote{卡点：“$V_i\ne0$ 当且仅当 $i$ 是指数”被称作众所周知，未给证明；$D_{2n}$ 的重数例外也显得突兀。背景：$\operatorname{ad}\check\rho$ 给出主分次，$\operatorname{ad}p_{-1}$ 把 $\mathfrak{b}_{i+1}$ 单射进 $\mathfrak{b}_i$，故每一步的余维数由该分次的维数决定；Kostant 的指数定理说这些余维数正是指数的重数，而指数等于 $U(\mathfrak{g})$ 中心生成元的次数减 $1$。结论：$E=\{d_1,\dots,d_\ell\}$ 是计重数的指数集合，$D_{2n}$ 中有一个指数出现两次，对应 $U(\mathfrak{g})$ 中心的两个同次生成元。来源：原书 p.112；B. Kostant [1963]（本书文献表）。}

\begin{lemma}[Drinfeld and Sokolov 1985]
    $N[[t]]$ 在 $\widetilde{\Op}_G(D)$ 上的作用是自由的。
\end{lemma}

\begin{proof}
我们断言，$\widetilde{\Op}_G(D)$ 中每个元素 $\partial_t+p_{-1}+\mathbf{v}(t)$ 都可唯一地写成
\begin{equation}
    \partial_t+p_{-1}+\mathbf{v}(t)=\exp(\operatorname{ad}U)\cdot(\partial_t+p_{-1}+\mathbf{c}(t)),
    \tag{4.2-4}
\end{equation}
其中 $U\in\mathfrak{n}[[t]]$，$\mathbf{c}(t)\in V[[t]]$。为此按主分次分解：$U=\sum_{j\ge0}U_j$，$\mathbf{v}(t)=\sum_{j\ge0}\mathbf{v}_j(t)$，$\mathbf{c}(t)=\sum_{j\in E}\mathbf{c}_j(t)$。比较 (4.2-4) 两边次数为 $j$ 的齐次分量，得到 $\mathbf{c}_i+[U_{i+1},p_{-1}]$ 可用 $\mathbf{v}_i$、$\mathbf{c}_j\,(j<i)$ 与 $U_j\,(j\le i)$ 表示。由 $\operatorname{ad}p_{-1}$ 的单射性便可唯一确定 $\mathbf{c}_i$ 与 $U_{i+1}$。于是满足 (4.2-4) 的 $U$ 与 $\mathbf{c}$ 可由归纳法唯一求得，引理得证。
\end{proof}

% [Tutor Note: 这个证明是“主 sl_2 三元组”理论的标准逐次消去（Kostant 的指数定理）：$\operatorname{ad}p_{-1}$ 在主分次下把 $\mathfrak{b}_{i+1}$ 单射进 $\mathfrak{b}_i$，故每一步只留下 $V_i$ 无法消去，而 $V_i\ne0$ 恰对应指数。原文说“众所周知”的部分即 $V_0=0$ 与 $\dim V_i=$ 重数，见 [Kostant 1963] 或 [Drinfeld and Sokolov 1985] §1。]

横截子空间 $V=\bigoplus_{i\in E}V_i$ 有一个特殊选取。即存在唯一的 $p_1\in\mathfrak{n}$，使得 $\{p_{-1},2\check{\rho},p_1\}$ 是 $\mathfrak{sl}_2$-三元组。这意味着它们满足与 $\mathfrak{sl}_2$ 的生成元 $\{e,h,f\}$ 相同的对易关系（见 §2.1.1）。

例如，对 $\mathfrak{g}=\mathfrak{sl}_n$（秩 $\ell=n-1$）有
\[
    p_{-1}=\begin{pmatrix}
        0&&&&\\
        1&0&&&\\
        &1&0&&\\
        &&\ddots&\ddots&\\
        &&&1&0
    \end{pmatrix},
\]
\[
    p_0=\begin{pmatrix}
        n-1&&&&\\
        &n-3&&&\\
        &&n-5&&\\
        &&&\ddots&\\
        &&&&-n+1
    \end{pmatrix},
\]
\[
    p_1=\begin{pmatrix}
        0&1(n-1)&&&\\
        &0&2(n-2)&&\\
        &&\ddots&\ddots&\\
        &&&0&(n-1)1\\
        &&&&0
    \end{pmatrix}.
\]
一般地，$p_1=\sum_{i=1}^{\ell}m_ie_i$，其中 $e_i$ 是 $\mathfrak{n}_+$ 的生成元，$m_i$ 是某些系数，由 $\{p_{-1},2\check{\rho},p_1\}$ 是 $\mathfrak{sl}_2$-三元组这一条件唯一确定。

设 $V^{\mathrm{can}}=\bigoplus_{i\in E}V_i^{\mathrm{can}}$ 是 $\mathfrak{n}$ 中 $\operatorname{ad}p_1$-不变元组成的空间。则 $p_1$ 张成 $V_1^{\mathrm{can}}$。设 $p_j$ 是 $V_{d_j}^{\mathrm{can}}$ 的线性生成元。若 $d_j$ 的重数大于 $1$，则在 $V_{d_j}^{\mathrm{can}}$ 中选取线性无关的向量。

$\mathfrak{g}=\mathfrak{sl}_n$ 时，可取矩阵 $p_1^j$（矩阵乘积）作为元素 $p_j,\,j=1,\dots,n-1$。

由引理 4.2.2，每个 $G$-oper 都可由唯一的算子 $\nabla=\partial_t+p_{-1}+\mathbf{v}(t)$ 代表，其中 $\mathbf{v}(t)\in V^{\mathrm{can}}[[t]]$，于是可写
\[
    \mathbf{v}(t)=\sum_{j=1}^{\ell}v_j(t)\cdot p_j,\qquad v_j(t)\in\mathbb{C}[[t]].
\]

下面确定坐标变换群如何作用在标准代表元上。

设 $t=\varphi(s)$，其中 $s$ 是 $D_x$ 上的另一坐标。关于新坐标 $s$，$\nabla$ 变为 $\partial_s+\widetilde{\mathbf{v}}(s)$，其中 $\widetilde{\mathbf{v}}(s)$ 按公式 (4.2-3) 由 $\mathbf{v}(t)$ 与 $\varphi(s)$ 表出。由引理 4.2.2，存在唯一的算子 $\partial_s+p_{-1}+\overline{\mathbf{v}}(s)$（其中 $\overline{\mathbf{v}}(s)\in V^{\mathrm{can}}[[s]]$）与 $g\in B[[s]]$，使得
\begin{equation}
    \partial_s+p_{-1}+\overline{\mathbf{v}}(s)=g\cdot\left(\partial_s+\widetilde{\mathbf{v}}(s)\right).
    \tag{4.2-5}
\end{equation}
直接计算可得
\begin{equation}
    g=\exp\left(\frac{1}{2}\frac{\varphi''}{\varphi'}\cdot p_1\right)\check{\rho}(\varphi'),
    \tag{4.2-6}
\end{equation}
\begin{equation}
    \overline{v}_1(s)=v_1(\varphi(s))(\varphi')^2-\frac{1}{2}\{\varphi,s\},
    \tag{4.2-7}
\end{equation}
\begin{equation}
    \overline{v}_j(s)=v_j(\varphi(s))(\varphi')^{d_j+1},\qquad j>1,
    \tag{4.2-8}
\end{equation}
其中 $\{\varphi,s\}$ 是 Schwarz 导数 (3.5-12)。

% [Tutor Note: 这三条公式必须一起核对符号：(4.2-6) 指数中 $\frac12\frac{\varphi''}{\varphi'}$ 的 $1/2$ 与 (4.2-7) 中 Schwarz 导数的 $-\frac12$ 严格配对；若符号反了，$v_1$ 的变换就不再是射影联络的 (4.1-2)。另注意 (4.2-7) 中 $(\varphi')^2$ 对应 $d_1=1$，(4.2-8) 中 $(\varphi')^{d_j+1}$ 说明 $v_j$ 是 $(d_j+1)$-微分。]

公式 (4.2-6) 可用来描述 $B$-丛 $\mathcal{F}_B$。即在 §4.1.3 中我们认出了所有 $PGL_2$-opers 底下的 $PGL_2$-丛 $\mathcal{F}_{PGL_2}$ 及其 $B_{PGL_2}$-约化 $\mathcal{F}_{B_{PGL_2}}$。设 $\iota:PGL_2\hookrightarrow G$ 是对应于我们一直使用的 $\mathfrak{g}$ 的主 $\mathfrak{sl}_2$ 子代数的主嵌入（回忆按假设 $G$ 是伴随型）。用 $\iota_B$ 记相应的嵌入 $B_{PGL_2}\hookrightarrow B$。则按公式 (4.2-6)，所有 $G$-opers 底下的 $G$-丛 $\mathcal{F}$ 及其 $B$-约化 $\mathcal{F}_B$，分别同构于由 $\mathcal{F}_{PGL_2}$ 与 $\mathcal{F}_{B_{PGL_2}}$ 沿 $\iota$ 与 $\iota_B$ 诱导出的丛。

公式 (4.2-7) 与 (4.2-8) 表明：在坐标变换下，$v_1$ 按射影联络变换，而 $v_j\,(j>1)$ 按 $(d_j+1)$-微分变换。于是在 $D_x$ 上得到同构
\begin{equation}
    \Op_G(D_x)\simeq\mathcal{P}roj(D_x)\times\bigoplus_{j=2}^{\ell}\Omega_{\mathcal{O}_x}^{\otimes(d_j+1)},
    \tag{4.2-9}
\end{equation}
其中 $\Omega_{\mathcal{O}_x}^{\otimes n}$ 是 $D_x$ 上 $n$-微分的空间，$\mathcal{P}roj(D_x)$ 是 $D_x$ 上射影联络的 $\Omega_{\mathcal{O}_x}^{\otimes2}$-torsor。

此分析逐字照搬到去心圆盘 $D_x^\times$ 的情形。特别地，对 $D_x$ 上坐标 $t$ 的每个选取，得到 $\Op_G(D_x^\times)$ 与算子 $\nabla=\partial_t+p_{-1}+\mathbf{v}(t)$（此时 $\mathbf{v}(t)\in V^{\mathrm{can}}((t))$）的空间之间的一个等同。坐标变换的作用由同样的公式给出，于是得到同构 (4.2-9) 的类似物：
\begin{equation}
    \Op_G(D_x^\times)\simeq\mathcal{P}roj(D_x^\times)\times\bigoplus_{j=2}^{\ell}\Omega_{\mathcal{K}_x}^{\otimes(d_j+1)}.
    \tag{4.2-10}
\end{equation}
自然嵌入 $\Op_G(D_x)\hookrightarrow\Op_G(D_x^\times)$ 与 (4.2-9)、(4.2-10) 右端的自然嵌入相容。

同样地，对一般光滑曲线 $X$ 得到等同
\[
    \Op_G(X)\simeq\mathcal{P}roj(X)\times\bigoplus_{j=2}^{\ell}\Gamma\left(X,\Omega_X^{\otimes(d_j+1)}\right).
\]

这个“标准形”到底有多典范？设除 $\mathfrak{g}$ 中 Borel 子代数 $\mathfrak{b}$ 之外不作任何选取。标准形的定义涉及 $\mathfrak{sl}_2$-三元组 $\{p_{-1},2\check{\rho},p_1\}$ 的选取。但选取这样的 $\mathfrak{sl}_2$-三元组，等价于选取投影 $\mathfrak{b}/\mathfrak{n}\to\mathfrak{h}$ 的分裂 $\mathfrak{h}\hookrightarrow\mathfrak{b}$，并选取相应一维子空间 $\mathfrak{n}_{-\alpha_i}$ 的生成元。容易看出，群 $B$ 单传递地作用在这些选取之集上，从而作用在 $\mathfrak{sl}_2$-三元组之集上。

给定 $\mathfrak{sl}_2$-三元组 $\{p_{-1},2\check{\rho},p_1\}$，我们还选取了 $V^{\mathrm{can}}$ 的基元素 $p_i,\,i\in E$，它们关于 $\check{\rho}$ 定义的分次是齐次的。这使我们能作出等同 (4.2-9)。然而即使 $d_i$ 的重数为 $1$，这些元素也只是在相差一个标量的意义下确定。因此若要真正典范的形式，应把 (4.2-9) 中用的向量空间 $\bigoplus_{i=2}^{\ell}\Omega_{\mathcal{O}_x}^{d_i}\cdot p_i$ 换成空间
\[
    V_{>1,x}^{\mathrm{can}}=\Gamma\left(D_x,\Omega_X^\times\times_{\mathbb{C}^\times}V_{>1}^{\mathrm{can}}(1)\right),
\]
其中 $V_{>1}^{\mathrm{can}}(1)$ 是向量空间 $\bigoplus_{i>0}V_i^{\mathrm{can}}$ 把分次平移 $1$，它不再依赖基的选取。于是得到更典范的等同
\begin{equation}
    \Op_G(D_x)\simeq\mathcal{P}roj(D_x)\times V_{>1,x}^{\mathrm{can}},
    \tag{4.2-11}
\end{equation}
以及 $D_x^\times$ 与更一般曲线 $X$ 的类似物；后者现在只依赖 $\mathfrak{sl}_2$-三元组的选取。特别地，由于 $\mathcal{P}roj(D_x)$ 是 $\Omega_{\mathcal{O}_x}^{\otimes2}$-torsor，通过把 $\Omega_{\mathcal{O}_x}^{\otimes2}=\Omega_{\mathcal{O}_x}^2\cdot p_1$ 等同，我们得到 $\Op_G(D_x)$ 是 $V_x^{\mathrm{can}}$-torsor，其中
\[
    V_x^{\mathrm{can}}=\Gamma\left(D_x,\Omega_X^\times\times_{\mathbb{C}^\times}V^{\mathrm{can}}(1)\right).
\]

现设另有 $\mathfrak{sl}_2$-三元组 $\{\widetilde{p}_{-1},2\check{\rho},\widetilde{p}_1\}$ 及相应的空间 $\widetilde{V}^{\mathrm{can}}$。两者都由原来的三元组经唯一确定的元素 $b\in B$ 的伴随作用得到。给定一个 oper $\rho$，得到它的两种不同标准形：
\[
    \nabla=\partial_t+p_{-1}+v_1(t)p_1+\mathbf{v}_{>1}(t),\qquad
    \widetilde{\nabla}=\partial_t+\widetilde{p}_{-1}+\widetilde{v}_1(t)\widetilde{p}_1+\widetilde{\mathbf{v}}_{>1}(t),
\]
其中 $\mathbf{v}(t)\in V_{>1}^{\mathrm{can}}$，$\widetilde{\mathbf{v}}(t)\in\widetilde{V}_{>1}^{\mathrm{can}}$。但由构造可知 $\widetilde{\nabla}=b\nabla b^{-1}$，这意味着 $\widetilde{v}_1(t)=v_1(t)$，且
\[
    \widetilde{\mathbf{v}}_{>1}(t)=b\mathbf{v}_{>1}(t)b^{-1}.
\]
由于所有可能的子空间 $V_{>1}^{\mathrm{can}}$ 都按此方式彼此典范等同，可把它们全部等同于一个唯一的“抽象”分次向量空间 $V_{<1}^{\mathrm{abs}}$（但注意 $V_{<1}^{\mathrm{abs}}$ 不能典范地等同于 $\mathfrak{b}$ 的子空间）。按前面的方式定义它的扭 $V_{>1,x}^{\mathrm{abs}}$。则 opers 的“真正”标准形是同构
\begin{equation}
    \Op_G(D_x)\simeq\mathcal{P}roj(D_x)\times V_{>1,x}^{\mathrm{abs}},
    \tag{4.2-12}
\end{equation}
它不依赖任何选取（$D_x^\times$ 与更一般曲线 $X$ 的情形类似）。特别地，如上可得 $\Op_G(D_x)$ 是 $V_x^{\mathrm{abs}}$-torsor，其中 $V_x^{\mathrm{abs}}$ 是按前述方式定义的抽象分次空间 $V^{\mathrm{abs}}$ 的扭。

% [Tutor Note: 原文在这里把抽象分次空间记作 $V^{\mathrm{abs}}_{<1}$（连写两次），但紧接着定义其扭时写的是 $V^{\mathrm{abs}}_{>1,x}$，且上一段刚用过 $V^{\mathrm{can}}_{>1}$。按上下文，抽象空间应是 $V^{\mathrm{abs}}_{>1}$（下标 $<1$ 是笔误，很可能由 $V^{\mathrm{can}}_{>1}$ 的排版串行所致）。译文保留原文写法，请读者按 $>1$ 理解。]

\subsection{$\mathfrak{sl}_n$ 的另一种代表元选取}

$\mathfrak{g}=\mathfrak{sl}_n$ 时还有另一种在应用中很有用的 oper 规范类代表元选取。即取横截子空间 $V=\bigoplus_{i=1}^{n-1}V_i$ 为只第一行有非零元的无迹矩阵空间（于是其第 $i$ 个分量 $V_i$ 是矩阵 $E_{1,i+1}$ 张成的空间）。相应的 $PGL_n$-opers 代表元具有形式
\begin{equation}
    \partial_t+\begin{pmatrix}
        0&u_1&u_2&\cdots&u_{n-1}\\
        1&0&0&\cdots&0\\
        0&1&0&\cdots&0\\
        \vdots&\ddots&\ddots&\ddots&\vdots\\
        0&0&\cdots&1&0
    \end{pmatrix}.
    \tag{4.2-13}
\end{equation}
与 $\mathfrak{sl}_2$ 的情形一样可以证明，该空间与 $n$ 阶微分算子
\begin{equation}
    \partial_t^n-u_1(t)\partial_t^{n-2}+\cdots+u_{n-2}(t)\partial_t-(-1)^nu_{n-1}(t)
    \tag{4.2-14}
\end{equation}
（作用 $\Omega^{-(n-1)/2}\to\Omega^{(n+1)/2}$）的空间重合。它们的主符号为 $1$，次主符号为 $0$。

这种形式的好处是非常显式，给出 $PGL_n$-opers 用标量微分算子的具体实现。其他经典李群也有类似的实现（见 [Drinfeld and Sokolov 1985; Beilinson and Drinfeld 2005]）。缺点是它不够典范，因为“第一行”子空间并非典范定义。此外，算子 (4.2-14) 的系数 $u_i(t)$ 在坐标变换下变换规律复杂，不同于上一节标准形的系数 $v_i(t)$。\transnote{卡点：(4.2-14) 为何作用 $\Omega^{-(n-1)/2}\to\Omega^{(n+1)/2}$，末项符号 $(-1)^n$ 又从何而来。背景：这是 $\mathfrak{sl}_n$ 的 Drinfeld–Sokolov 形式：$n$ 阶标量算子的主符号为 $1$，$\partial_t$ 每作用一次把 $\Omega^s$ 送到 $\Omega^{s+1}$，故 $s=-(n-1)/2$ 时像落在 $\Omega^{(n+1)/2}$。结论：$(-1)^n$ 是把 $u_{n-1}$ 与 $\mathfrak{sl}_2$-三元组的规范化相匹配的约定，换文献时极易差一个符号。来源：原书 p.119–120 与 [Drinfeld and Sokolov 1985]（本书文献表）。}

% [Tutor Note: (4.2-14) 就是经典的 Drinfeld–Sokolov 形式：$n$ 阶标量算子与 $PGL_n$-oper 的对应。核对时注意最后一项的符号 $(-1)^n$，它来自把 $u_{n-1}$ 放进伴随表示所需的 $\mathfrak{sl}_2$-三元组规范化；不同文献把 $u_{n-1}$ 定义成 $(-1)^{n-1}u_{n-1}$，抄写时极易差一个符号。]

\section{任意仿射 Kac--Moody 代数的中心}

现在陈述本书的主要结果之一：顶点代数 $V_{\kappa_c}(\mathfrak{g})$ 的中心 $\mathfrak{z}(\hat{\mathfrak{g}})$ 的坐标无关描述。这个描述由定理 4.3.1 与 4.3.2 给出，推广了定理 3.5.2 中 $\mathfrak{z}(\widehat{\mathfrak{sl}_2})$ 的描述。我们将用它来描述临界水平下完备泛包络代数 $\widetilde{U}_{\kappa_c}(\hat{\mathfrak{g}})$ 的中心 $Z(\hat{\mathfrak{g}})$（定理 4.3.6）。最后在命题 4.3.9 中证明：离开临界水平时，$\widetilde{U}_\kappa(\hat{\mathfrak{g}})$ 的中心是平凡的。

\begin{keypoints}
    \item \textbf{核心思想：}顶点代数的中心由 Langlands 对偶群 ${}^LG$ 的 opers 空间上的函数给出（Feigin–Frenkel 定理），这是本章的主定理。
    \item \textbf{数学工具：}形式特征标、PBW 滤过与分次伴随代数、反向极限、$\mathfrak{g}[[t]]$-不变函数。
    \item \textbf{关键结论：}$\mathfrak{z}(\hat{\mathfrak{g}})\simeq\operatorname{Fun}\Op_{{}^LG}(D)$；$Z(\hat{\mathfrak{g}})\simeq\operatorname{Fun}\Op_{{}^LG}(D^\times)$；$G((t))$ 在中心上作用平凡；$\kappa\ne\kappa_c$ 时中心为标量。
    \item \textbf{记号提醒：}$\mathfrak{z}(\hat{\mathfrak{g}})$ 是顶点代数中心，$Z(\hat{\mathfrak{g}})$ 是完备泛包络代数中心，二者由 $\widetilde{U}$ 函子相联系。
\end{keypoints}

\subsection{顶点代数的中心}

回忆 §3.5.2：对曲线 $X$ 的任一光滑点 $x$，有联系于 $x$ 的仿射 Kac--Moody 代数 $\hat{\mathfrak{g}}_x$ 以及由公式 (3.5-4) 定义的代数 $\mathfrak{z}(\hat{\mathfrak{g}})_x$。

设 ${}^LG$ 是联系于 $\mathfrak{g}$ 的 Langlands 对偶李代数 ${}^L\mathfrak{g}$ 的伴随型单李群。按定义，${}^L\mathfrak{g}$ 是 Cartan 矩阵为 $\mathfrak{g}$ 的 Cartan 矩阵之转置的李代数（见附录 A.3）。于是 ${}^LG$ 是群 $G$ 的 Langlands 对偶群（按 §1.1.5 的定义），这里 $G$ 是以 $\mathfrak{g}$ 为李代数的连通单连通单李群。

以下定理是本书的核心结果，由 B. Feigin 与我本人证明 [1992]（另见 [Frenkel 2005b]）。其详细证明将在后面几章给出。

\begin{theorem}
    代数 $\mathfrak{z}(\hat{\mathfrak{g}})_x$ 自然同构于圆盘 $D_x$ 上 ${}^LG$-opers 的空间 $\Op_{{}^LG}(D_x)$ 上的函数代数。
\end{theorem}

等价地，可考虑对应于标准圆盘 $D=\operatorname{Spec}\mathcal{O}$（$\mathcal{O}=\mathbb{C}[[t]]$）的代数 $\mathfrak{z}(\hat{\mathfrak{g}})$，但记住群 $\operatorname{Aut}\mathcal{O}$ 与李代数 $\operatorname{Der}\mathcal{O}$ 的作用。它们作用在 $\mathfrak{z}(\hat{\mathfrak{g}})\subset V_{\kappa_c}(\mathfrak{g})$ 与空间 $\Op_{{}^LG}(D)$（从而作用在函数代数 $\operatorname{Fun}\Op_{{}^LG}(D)$）上。\transnote{卡点：定理 4.3.1 与定理 3.5.2 的关系，以及临界水平 $\kappa_c$ 与对偶群 ${}^LG$ 如何进入。背景：$\kappa_c=-h^\vee\kappa_0$ 是使中心最大的水平（§3.5.2），Feigin–Frenkel 定理断言 $\mathfrak{z}(\hat{\mathfrak{g}})\cong\operatorname{Fun}\Op_{{}^LG}(D)$；当 $\mathfrak{g}=\mathfrak{sl}_2$ 时 ${}^LG=PGL_2$，oper 就是射影联络，$\operatorname{Fun}\Op\cong\mathbb{C}[v_{1,n}]_{n<0}$，正回到定理 3.5.2。结论：这是“中心 $=$ 对偶群 oper 空间上的函数”的精确表述，其证明占本书主要篇幅。来源：原书 p.117 与 [Feigin and Frenkel 1992]（本书文献表）。}

\begin{theorem}
    中心 $\mathfrak{z}(\hat{\mathfrak{g}})$ 与代数 $\operatorname{Fun}\Op_{{}^LG}(D)$ 同构，且该同构是 $(\operatorname{Der}\mathcal{O},\operatorname{Aut}\mathcal{O})$-等变的。
\end{theorem}

% [Tutor Note: 定理 4.3.1 中的群是 ${}^LG$ 而不是 $G$——这正是 §4.2 引言所说的“转折”。若把 $G$ 与 ${}^LG$ 弄混，结论在 $\mathfrak{g}$ 非自对偶时（如 $B_\ell$ 与 $C_\ell$）会明显错误。另注意定理中的 $G$ 取单连通，而 ${}^LG$ 取伴随型：两者的中心不同，这解释了为什么 ${}^LG$ 的 B-约化与余权格能如此干净地对应 opers。]

利用同构 (4.2-9)，得到同构
\[
    \Op_{{}^LG}(D)\simeq\mathbb{C}[[t]]^{\oplus\ell}
\]
（当然它依赖坐标 $t$）。因此每个 oper 由 $\ell$ 元形式 Taylor 级数 $(v_1(t),\dots,v_\ell(t))$ 代表。写
\[
    v_i(t)=\sum_{n<0}v_{i,n}t^{-n-1}.
\]
于是得到同构
\begin{equation}
    \operatorname{Fun}\Op_{{}^LG}(D)\simeq\mathbb{C}[v_{i,n}]_{i=1,\dots,\ell;\,n<0}.
    \tag{4.3-1}
\end{equation}

设 $S_i,\,i=1,\dots,\ell$，是 $\mathfrak{z}(\hat{\mathfrak{g}})$ 中在定理 4.3.2 的同构下对应于 $v_{i,-1}\in\operatorname{Fun}\Op_{{}^LG}(D)$ 的元素。下面利用该同构关于算子 $L_{-1}=-\partial_t\in\operatorname{Der}\mathcal{O}$ 的作用的等变性。在 $\operatorname{Fun}\Op_{{}^LG}(D)$ 一侧显然有 $\frac{1}{m!}(-\partial_t)^mv_{i,-1}=v_{i,-m-1}$。于是得到：在定理 4.3.2 的同构下，生成元 $v_{i,-m-1}$ 映到
\[
    \frac{1}{m!}L_{-1}^mS_i=\frac{1}{m!}T^mS_i,
\]
因为 $L_{-1}$ 与顶点代数 $V_{\kappa_c}(\mathfrak{g})$ 上的平移算子 $T$ 一致。这意味着
\[
    v_{i,-m-1}\mapsto S_{i,(-m-1)}|0\rangle,
\]
其中 $S_{i,(n)}$ 是顶点算子
\[
    Y(S_i,z)=\sum_{n<0}S_{i,(n)}z^{-n-1}
\]
的 Fourier 系数。事实上，回忆 $\mathfrak{z}(\hat{\mathfrak{g}})$ 是交换顶点代数，故可用公式 (2.2-1)。注意算子 $S_{i,(n)}\,(n\ge0)$ 零化真空向量 $|0\rangle$，且由于它们与 $\hat{\mathfrak{g}}_{\kappa_c}$ 的作用交换，其在 $V_{\kappa_c}(\mathfrak{g})$ 上的作用恒等于零。

此外，$v_{i,n}$ 的任一多项式在定理 4.3.2 的同构下映到 $S_{i,(n)}$ 的相应多项式作用于真空向量 $|0\rangle\in V_{\kappa_c}(\mathfrak{g})$。于是
\begin{equation}
    \mathfrak{z}(\hat{\mathfrak{g}})=\mathbb{C}[S_{i,(n)}]_{i=1,\dots,\ell;\,n<0}|0\rangle.
    \tag{4.3-2}
\end{equation}
这个等同特别意味着：引理 3.3.1 中的映射 $\operatorname{gr}\mathfrak{z}(\hat{\mathfrak{g}})\hookrightarrow\operatorname{Inv}\mathfrak{g}^*[[t]]$ 是同构。\footnote{我们指出，这其实是下面给出的定理 4.3.2 证明中的一步。}

为看出这点，回忆定理 3.4.2 中代数 $\operatorname{Inv}\mathfrak{g}^*[[t]]$ 的描述：
\[
    \operatorname{Inv}\mathfrak{g}^*[[t]]=\mathbb{C}[\overline{P}_{i,n}]_{i=1,\dots,\ell;\,n<0},
\]
其中 $\overline{P}_i,\,i=1,\dots,\ell$，是代数 $\operatorname{Inv}\mathfrak{g}^*$ 的次数为 $d_i+1$ 的齐次生成元。要比较 $\operatorname{gr}\mathfrak{z}(\hat{\mathfrak{g}})$ 与 $\operatorname{Inv}\mathfrak{g}^*[[t]]$ 这两个空间，我们计算它们的形式特征标，即关于分次算子 $L_0=-t\partial_t$ 的分次维数的生成函数。

若 $V$ 是分次分量有限维的 $\mathbb{Z}$-分次向量空间 $V=\bigoplus_{n\in\mathbb{Z}}V_n$，定义其形式特征标为
\[
    \operatorname{ch}V=\sum_{n\in\mathbb{Z}}\dim V_nq^n.
\]

现在，由 §3.3.4 中给出的 $\operatorname{Inv}\mathfrak{g}^*[[t]]$ 的生成元 $\overline{P}_{i,n}$ 的定义可知 $\deg\overline{P}_{i,m}=d_i-m$。因此
\begin{equation}
    \operatorname{ch}\operatorname{Inv}\mathfrak{g}^*[[t]]=\prod_{i=1}^{\ell}\prod_{n_i\ge d_i+1}(1-q^{n_i})^{-1}.
    \tag{4.3-3}
\end{equation}
另一方面，由公式 (4.2-9) 给出的 $\Op_{{}^LG}(D)$ 的描述，可得 $-t\partial_t\in\operatorname{Der}\mathcal{O}$ 对生成元 $v_{i,m}$ 的作用。算出 $\deg v_{i,m}=d_i-m$。故 $\operatorname{Fun}\Op_{{}^LG}(D)$ 的特征标（从而 $\mathfrak{z}(\hat{\mathfrak{g}})$ 的特征标）等于公式 (4.3-3) 的右端。既然我们有特征标相同的两个分次空间之间的单射 $\operatorname{gr}\mathfrak{z}(\hat{\mathfrak{g}})\hookrightarrow\operatorname{Inv}\mathfrak{g}^*[[t]]$，便得到：

\begin{proposition}
    自然嵌入 $\operatorname{gr}\mathfrak{z}(\hat{\mathfrak{g}})\hookrightarrow\operatorname{Inv}\mathfrak{g}^*[[t]]$ 是同构。
\end{proposition}

换言之，中心 $\mathfrak{z}(\hat{\mathfrak{g}})$ 是“尽可能大的”。这意味着 $\operatorname{Inv}\mathfrak{g}^*[[t]]$ 的生成元 $\overline{P}_{i,-1}$ 可被“提升”为 $\mathfrak{g}[[t]]$-不变向量 $S_i\in\mathfrak{z}(\hat{\mathfrak{g}})$，$i=1,\dots,\ell$。换言之，$\overline{P}_{i,-1}$ 是 $S_i$ 的符号。由此还得到：对所有 $n<0$，$\overline{P}_{i,n}$ 是 $S_{i,(n)}|0\rangle\in\mathfrak{z}(\hat{\mathfrak{g}})$ 的符号。\transnote{卡点：为何“形式特征标相同 + 分次单射”就能推出同构。背景：两边都是分次分量有限维的 $\mathbb{Z}$-分次空间，单射逐分量给出 $\dim(\operatorname{gr}\mathfrak{z})_n\le\dim(\operatorname{Inv})_n$。结论：形式幂级数相等意味着系数逐项相等，故每个分量的维数相等，单射自动是同构；这一步不需要有限生成或无挠假设。来源：依据原书 p.119–120 的证明。}

% [Tutor Note: (4.3-3) 中每个指数的重数在 $D_{2n}$、$d_n=2n$ 时为 2，故该公式在那一情形要把因子重复两次；这与 §4.2.4 中“重数等于 2”的附注是同一件事。另注意特征标相等只给“维数逐项相等”，把单射升级为同构靠的是分次映射保持分次，这一点原文略去未说。]

特别地，$\overline{P}_{1,-1}$ 是二次 Casimir 生成元，由公式
\[
    \overline{P}_{1,-1}=\frac{1}{2}J^a_{-1}J_{a,-1}
\]
给出（相差非零标量）。它到 $\mathfrak{z}(\hat{\mathfrak{g}})$ 的提升就是公式 (3.1-1) 给出的 Segal--Sugawara 向量 $S_1$。因此生成元 $S_{1,(n)}\,(n<0)$ 是 Segal--Sugawara 算子 $S_{n-1}$。

自然要问：能否把 $\operatorname{Inv}\mathfrak{g}^*[[t]]$ 的其他生成元 $\overline{P}_{i,-1}\,(i=1,\dots,\ell)$ 显式地提升到 $\mathfrak{z}(\hat{\mathfrak{g}})$？迄今答案只在特殊情形已知。对 $\mathfrak{g}=\mathfrak{sl}_n$，可用 R. Goodman 与 N. Wallach [1989] 的结果构造 $\overline{P}_{i,-1}$ 的显式公式，那里依赖一些复杂的不变量理论计算。\footnote{最近，[Chervov and Talalaev 2006] 对 $\mathfrak{g}=\mathfrak{sl}_n$ 得到了另一类优雅的公式。} 对经典类型 $\mathfrak{g}$（$A_\ell,B_\ell,C_\ell$），T. Hayashi 在 [1988] 中提出另一方法：他显式构造了 Segal--Sugawara 算子之后的接下来几个中心元，$A_\ell$ 情形次数为 $3$，$B_\ell$ 与 $C_\ell$ 情形次数为 $4$。随后他利用下文 §8.3.1 讨论的 $\widetilde{U}_{\kappa_c}(\hat{\mathfrak{g}})$ 中心上的 Poisson 结构（Hayashi 本质上正是为此引入该结构）生成更高次数的中心元。遗憾的是，当 $\mathfrak{g}=D_\ell$（Pfaffian 即 Pfaff 行列式，似乎难以触及）以及 $\mathfrak{g}$ 为例外李代数时，此方法不能给出整个中心。对 $A_\ell,B_\ell,C_\ell$ 型 $\mathfrak{g}$，这样可证明同构 (4.3-2)，但不能证明把中心 $\mathfrak{z}(\hat{\mathfrak{g}})$ 与代数 $\operatorname{Fun}\Op_{{}^LG}(D)$ 等同起来的定理 4.3.2。

B. Feigin 与我本人 [1992; 2005b] 给出的定理 4.3.2 的证明（本书所呈现的）不依赖中心生成元的显式公式。取而代之，我们把 $\mathfrak{z}(\hat{\mathfrak{g}})$ 与 $\operatorname{Fun}\Op_{{}^LG}(D)$ 都实现为另一个代数的子代数，即圆盘上某个 ${}^LH$-丛的联络空间上的函数代数，其中 ${}^LH$ 是 ${}^LG$ 的 Cartan 子群。

\subsection{泛包络代数的中心}

现在用上述关于 $\mathfrak{z}(\hat{\mathfrak{g}})$ 中心的描述，来描述临界水平下完备泛包络代数 $\widetilde{U}_{\kappa_c}(\hat{\mathfrak{g}})$ 的中心 $Z(\hat{\mathfrak{g}})$（定义见 §2.1.2）。

设 $B\in\mathfrak{z}(\hat{\mathfrak{g}})\subset V_{\kappa_c}(\mathfrak{g})$。则 $\mathfrak{g}[[t]]\cdot B=0$，且由对易关系公式 (3.2-2) 可知，§3.2.1 中定义的所有元素 $B_{[k]}\in\widetilde{U}_{\kappa_c}(\hat{\mathfrak{g}})$ 都与整个仿射 Kac--Moody 代数 $\hat{\mathfrak{g}}_{\kappa_c}$ 交换。因此它们是 $\widetilde{U}_{\kappa_c}(\hat{\mathfrak{g}})$ 的中心元。此外，$\widetilde{U}(\mathfrak{z}(\hat{\mathfrak{g}}))\subset\widetilde{U}(V_{\kappa_c}(\mathfrak{g}))=\widetilde{U}_{\kappa_c}(\hat{\mathfrak{g}})$ 的任一元素也是中心元（这里用到 §3.2.3 定义的泛包络代数函子 $\widetilde{U}$ 与引理 3.2.4）。于是得到同态 $\widetilde{U}(\mathfrak{z}(\hat{\mathfrak{g}}))\to Z(\hat{\mathfrak{g}})$。

现在，代数 $\widetilde{U}(\mathfrak{z}(\hat{\mathfrak{g}}))$ 是多项式代数的完备化。事实上，回忆上节我们有元素 $S_i\,(i=1,\dots,\ell)$ 给出同构 (4.3-2)。设 $S_{i,[n]}\,(n\in\mathbb{Z})$ 是 $\widetilde{U}_{\kappa_c}(\hat{\mathfrak{g}})$ 中对应于 $S_i\in V_{\kappa_c}(\mathfrak{g})$ 的元素（见 §3.2.1）。由定义可知，拓扑代数 $\widetilde{U}(\mathfrak{z}(\hat{\mathfrak{g}}))$ 是多项式代数 $\mathbb{C}[S_{i,[n]}]_{i=1,\dots,\ell;\,n\in\mathbb{Z}}$ 在如下拓扑下的完备化：$0$ 的开邻域基由 $S_{i,[n]}\,(i=1,\dots,\ell;\,n>N+d_i+1)$（$N\in\mathbb{Z}_+$）生成的理想构成。此拓扑等价于把条件 $n>N+d_i+1$ 换成 $n>N$ 的拓扑。

于是完备多项式代数 $\widetilde{U}(\mathfrak{z}(\hat{\mathfrak{g}}))$ 映入 $\widetilde{U}_{\kappa_c}(\hat{\mathfrak{g}})$ 的中心 $Z(\hat{\mathfrak{g}})$。

\begin{proposition}
    此映射是同构，故 $Z(\hat{\mathfrak{g}})$ 等于 $\widetilde{U}(\mathfrak{z}(\hat{\mathfrak{g}}))$。
\end{proposition}

\begin{proof}
我们沿用 [Beilinson and Drinfeld 1997] 定理 3.7.7 的证明。

先描述 $Z(\hat{\mathfrak{g}})$ 的分次伴随代数。$U_{\kappa_c}(\hat{\mathfrak{g}})$ 上的 Poincaré--Birkhoff--Witt 滤过在其完备化 $\widetilde{U}_{\kappa_c}(\hat{\mathfrak{g}})$ 上诱导滤过。分次伴随代数 $\operatorname{gr}U_{\kappa_c}(\hat{\mathfrak{g}})$（回忆 §3.2.3 的定义）是代数 $\operatorname{Sym}\mathfrak{g}((t))=\bigoplus_{i\ge0}\operatorname{Sym}^i\mathfrak{g}((t))$。设 $\overline{I}_N$ 是 $\operatorname{Sym}\mathfrak{g}((t))$ 中由 $\mathfrak{g}\otimes t^N\mathbb{C}[[t]]$ 生成的理想。分次伴随代数 $\operatorname{gr}\widetilde{U}_{\kappa_c}(\hat{\mathfrak{g}})$ 是 $\operatorname{Sym}\mathfrak{g}((t))$ 的完备化 $\widetilde{\operatorname{Sym}}\mathfrak{g}((t))=\bigoplus_{i\ge0}\widetilde{\operatorname{Sym}}^i\mathfrak{g}((t))$，其中
\[
    \widetilde{\operatorname{Sym}}^i\mathfrak{g}((t))=\lim_{\longleftarrow}\operatorname{Sym}^i\mathfrak{g}((t))/\left(\overline{I}_N\cap\operatorname{Sym}^i\mathfrak{g}((t))\right).
\]

换言之，$\widetilde{\operatorname{Sym}}^i\mathfrak{g}((t))$ 是 $\mathfrak{g}((t))\simeq\mathfrak{g}^*((t))dt\simeq\mathfrak{g}^*((t))$ 上次数为 $i$ 的多项式函数空间 $\operatorname{Fun}^i(\mathfrak{g}^*\otimes t^{-N}\mathbb{C}[[t]])$。所谓这样的函数，是指 $\mathfrak{g}^*((t))$ 上的函数，使得对每个 $N\in\mathbb{Z}_+$，它到 $\mathfrak{g}^*\otimes t^{-N}\mathbb{C}[[t]]$ 的限制是次数 $i$ 的多项式函数（即对充分大的 $n$，它来自有限维向量空间 $\mathfrak{g}^*\otimes(t^{-N}\mathbb{C}[[t]]/t^n\mathbb{C}[[t]])$ 上多项式函数的拉回）。

李代数 $\mathfrak{g}((t))$ 经伴随作用自然作用在每个空间 $\widetilde{\operatorname{Sym}}^i\mathfrak{g}((t))$ 上，其李子代数 $\mathfrak{g}[[t]]$ 保持子空间 $\overline{I}_N\cap\operatorname{Sym}^i\mathfrak{g}((t))$，即那些在 $\mathfrak{g}\otimes t^N\mathbb{C}[[t]]$ 上取零值的函数。

设 $\widetilde{\operatorname{Inv}}\mathfrak{g}^*((t))=\bigoplus_{i\ge0}\widetilde{\operatorname{Inv}}^i\mathfrak{g}((t))$ 是 $\widetilde{\operatorname{Sym}}\mathfrak{g}((t))$ 中 $\mathfrak{g}((t))$-不变元组成的子代数。对每个 $N\in\mathbb{Z}_+$，取模 $\overline{I}_N$ 的商给出满同态
\[
    \widetilde{\operatorname{Sym}}^i\mathfrak{g}((t))\longrightarrow\operatorname{Fun}^i(\mathfrak{g}^*\otimes t^{-N}\mathbb{C}[[t]]).
\]
$\widetilde{\operatorname{Inv}}^i\mathfrak{g}^*((t))$ 在此同态下的像，包含在 $\operatorname{Fun}^i(\mathfrak{g}^*\otimes t^{-N}\mathbb{C}[[t]])$ 中 $\mathfrak{g}[[t]]$-不变元的空间内。

% [Tutor Note: 原文此处记号不统一：“设 $\widetilde{\mathrm{Inv}}\mathfrak{g}^*((t))=\bigoplus_{i\ge0}\widetilde{\mathrm{Inv}}^i\mathfrak{g}((t))$”，左端是 $\mathfrak{g}^*$ 而右端的每个分次分量写成 $\mathfrak{g}$；紧接着又写“$\widetilde{\mathrm{Inv}}^i\mathfrak{g}^*((t))$ 的像”。按上下文（不变元是 $\mathrm{Sym}\mathfrak{g}((t))$ 的对偶上的对象），各分量应统一为 $\mathfrak{g}^*((t))$。译文保留原文这一处的写法。]

但我们有自然同构
\[
    \operatorname{Fun}(\mathfrak{g}^*\otimes t^{-N}\mathbb{C}[[t]])\simeq\operatorname{Fun}(\mathfrak{g}^*[[t]])
\]
（乘以 $t^N$），它与 $\mathfrak{g}[[t]]$ 的作用交换。既然由定理 3.4.2 已知 $\mathfrak{g}^*[[t]]$ 上 $\mathfrak{g}[[t]]$-不变函数的代数，便可用此同构求出 $\mathfrak{g}^*\otimes t^{-N}\mathbb{C}[[t]]$ 上 $\mathfrak{g}[[t]]$-不变函数的代数。为描述它，设 $\overline{J}^a_n$ 是 $\mathfrak{g}^*((t))$ 上的线性泛函，按公式 (3.3-4) 定义。注意 $\overline{J}^a_n\,(n\ge N)$ 到 $\mathfrak{g}^*\otimes t^{-N}\mathbb{C}[[t]]$ 的限制等于 $0$。记
\[
    \overline{P}_i(z)=\overline{P}_i(\overline{J}^a(z))=\sum_{n\in\mathbb{Z}}\overline{P}_{i,n}z^{-n-1},
\]
其中
\[
    \overline{J}^a(z)=\sum_{n\in\mathbb{Z}}\overline{J}^a_nz^{-n-1}.
\]
于是由上述同构得知，$\mathfrak{g}^*\otimes t^{-N}\mathbb{C}[[t]]$ 上 $\mathfrak{g}[[t]]$-不变函数的代数是生成元为 $\overline{P}_{i,n}\,(n<(d_i+1)N)$ 的自由多项式代数。

然而可以显式验证，每个 $\overline{P}_{i,n}$ 不仅 $\mathfrak{g}[[t]]$-不变，而且 $\mathfrak{g}((t))$-不变。这意味着同态
\[
    \widetilde{\operatorname{Inv}}\mathfrak{g}^*((t))\longrightarrow\left(\operatorname{Fun}(\mathfrak{g}^*\otimes t^{-N}\mathbb{C}[[t]])\right)^{\mathfrak{g}[[t]]}
\]
对所有 $N\in\mathbb{Z}_+$ 满射。于是 $\widetilde{\operatorname{Inv}}\mathfrak{g}^*((t))$ 是代数 $(\operatorname{Fun}(\mathfrak{g}^*\otimes t^{-N}\mathbb{C}[[t]]))^{\mathfrak{g}[[t]]}$ 的反向极限。

因此 $\widetilde{\operatorname{Inv}}\mathfrak{g}^*((t))$ 是多项式代数
\[
    \mathbb{C}[\overline{P}_{i,n}]_{i=1,\dots,\ell;\,n\in\mathbb{Z}}
\]
的完备化，其拓扑的 $0$ 开邻域基由固定次数且属于由 $\overline{P}_{i,n}\,(n>N,\ N\in\mathbb{Z}_+)$ 生成的理想的多项式子空间构成（用最后这个条件等价于用条件 $n\ge(d_i+1)N$，后者来自上面的分析）。

现在可以证明命题了。回忆 $\widetilde{U}_{\kappa_c}(\hat{\mathfrak{g}})$ 定义为 $U_{\kappa_c}(\hat{\mathfrak{g}})/I_N$ 之商的反向极限，其中 $I_N$ 是由 $\mathfrak{g}\otimes t^N\mathbb{C}[[t]]$ 生成的左理想。因此
\[
    Z(\hat{\mathfrak{g}})=\lim_{\longleftarrow}Z(\hat{\mathfrak{g}})/Z(\hat{\mathfrak{g}})\cap I_N.
\]
我们有单射 $Z(\hat{\mathfrak{g}})/Z(\hat{\mathfrak{g}})\cap I_N\hookrightarrow(U_{\kappa_c}(\hat{\mathfrak{g}})/I_N)^{\mathfrak{g}[[t]]}$。此外还有单射
\[
    \operatorname{gr}\left((U_{\kappa_c}(\hat{\mathfrak{g}})/I_N)^{\mathfrak{g}[[t]]}\right)\hookrightarrow\left(\operatorname{gr}U_{\kappa_c}(\hat{\mathfrak{g}})/\overline{I}_N\right)^{\mathfrak{g}[[t]]}=\left(\operatorname{Fun}(\mathfrak{g}^*\otimes t^{-N}\mathbb{C}[[t]])\right)^{\mathfrak{g}[[t]]}.
\]

已知 $\mathbb{C}[S_{i,[n]}]_{i=1,\dots,\ell;\,n\in\mathbb{Z}}$ 是 $Z(\hat{\mathfrak{g}})$ 的子代数。按构造，$S_{i,[n]}$ 的符号等于 $\overline{P}_{i,n}$。因此 $\mathbb{C}[S_{i,[n]}]_{i=1,\dots,\ell;\,n\in\mathbb{Z}}$ 在上述映射复合下的像就是整个 $(\operatorname{Fun}(\mathfrak{g}^*\otimes t^{-N}\mathbb{C}[[t]]))^{\mathfrak{g}[[t]]}$。这说明所有中间映射都是同构。特别地，对每个 $N\in\mathbb{Z}_+$ 有
\[
    Z(\hat{\mathfrak{g}})/Z(\hat{\mathfrak{g}})\cap I_N=\mathbb{C}[S_{i,[n_i]}]_{i=1,\dots,\ell;\,n_i\le N(d_i+1)}.
\]
因此中心 $Z(\mathfrak{g})$ 等于多项式代数 $\mathbb{C}[S_{i,[n]}]_{i=1,\dots,\ell;\,n\in\mathbb{Z}}$ 的完备化，其拓扑的 $0$ 开邻域基由 $S_{i,[n]}\,(i=1,\dots,\ell;\,n>N)$ 生成的理想构成。但这恰是 $\widetilde{U}(\mathfrak{z}(\hat{\mathfrak{g}}))$。证明完成。
\end{proof}

% [Tutor Note: $Z(\hat{\mathfrak{g}})/Z(\hat{\mathfrak{g}})\cap I_N$ 这条公式里的下标是 $n_i\le N(d_i+1)$，与前面定义的拓扑条件 $n>N$ 方向相反，这是完备化“截断”与“对偶”两套下标容易混淆之处：$I_N$ 截掉的是 $n>N(d_i+1)$ 的部分，所以商里剩下的是 $n_i\le N(d_i+1)$。另原文此句写作 $Z(\mathfrak{g})$（$Z$ 后无 hat），按上下文应为 $Z(\hat{\mathfrak{g}})$，$\mathfrak{g}$ 的中心在临界水平不可能是这个完备多项式代数。]

作为推论得到：$Z(\hat{\mathfrak{g}})$ 同构于 $\widetilde{U}(\operatorname{Fun}\Op_{{}^LG}(D))$。

\begin{lemma}
    代数 $\widetilde{U}(\operatorname{Fun}\Op_{{}^LG}(D))$ 典范同构于去心圆盘 $D^\times=\operatorname{Spec}\mathbb{C}((t))$ 上 ${}^LG$-opers 空间上函数构成的拓扑代数 $\operatorname{Fun}\Op_{{}^LG}(D^\times)$。
\end{lemma}

\begin{proof}
如 §4.3.1 所述，代数 $\operatorname{Fun}\Op_{{}^LG}(D)$ 同构于多项式代数 $\mathbb{C}[v_{i,n}]_{i=1,\dots,\ell;\,n<0}$，其中 $v_{i,n}$ 是 oper 联络
\begin{equation}
    \nabla=\partial_t+p_{-1}+\sum_{i=1}^{\ell}v_i(t)p_i
    \tag{4.3-4}
\end{equation}
的系数，$v_i(t)=\sum_{n<0}v_{i,n}t^{-n-1}$。函子 $\widetilde{U}$ 的构造推出：代数 $\widetilde{U}(\Op_{{}^LG}(D)_{\kappa_0}^{\vee})$ 是变量 $v_{i,n}\,(i=1,\dots,\ell;\,n\in\mathbb{Z})$ 的多项式代数在如下拓扑下的完备化：$0$ 开邻域基由 $v_{i,n}\,(n>N)$ 生成的理想构成。但这恰是形如 (4.3-4) 的 oper 空间上的函数代数，其中 $v_i(t)=\sum_{n\in\mathbb{Z}}v_{i,n}t^{-n-1}\in\mathbb{C}((t))$。\transnote{卡点：引理证明中的 $\widetilde U(\Op_{{}^LG}(D)^\vee_{\kappa_0})$ 下标 $\kappa_0$ 在此没有定义，且上一句刚把 $\operatorname{Fun}\Op_{{}^LG}(D)$ 等同于 $\mathbb{C}[v_{i,n}]$。背景：$\widetilde U$ 在此把变量 $v_{i,n}$ 的多项式代数按 $n>N$ 的拓扑完备化，上标 $\vee$ 表示取对偶。结论：读者可按 $\widetilde U(\Op_{{}^LG}(D)^\vee)$ 理解，$\kappa_0$ 系原书排版遗留，不影响本引理的论证。来源：原书 p.121–122（PDF p.135–136）。}
\end{proof}

% [Tutor Note: $\widetilde{U}(\Op_{{}^LG}(D)_{\kappa_0}^{\vee})$ 中的下标 $\kappa_0$ 在此处没有定义，且前一句刚把 $\mathrm{Fun}\,\Op_{{}^LG}(D)$ 等同于 $\mathbb{C}[v_{i,n}]$。按上下文，$\widetilde{U}$ 作用在 oper 空间的对偶（即系数空间 $v_{i,n}$）上，$\widetilde{U}$ 就是完备对称代数；上标 $\vee$ 正是取对偶，下标 $\kappa_0$ 疑为排版遗留。读者可按 $\widetilde{U}(\Op_{{}^LG}(D)^{\vee})$ 理解，这不影响本引理的论证。]

现在，定理 4.3.2、命题 4.3.4 与引理 4.3.5 蕴含：

\begin{theorem}
    中心 $Z(\hat{\mathfrak{g}})$ 与代数 $\operatorname{Fun}\Op_{{}^LG}(D^\times)$ 同构，且该同构是 $(\operatorname{Der}\mathcal{O},\operatorname{Aut}\mathcal{O})$-等变的。
\end{theorem}

后面将看到，此同构还满足其他一些条件，这些条件几乎唯一地确定它。

利用该定理，可描述李代数 $\hat{\mathfrak{g}}_{\kappa_c,x}$ 的泛包络代数 $\widetilde{U}_{\kappa_c}(\hat{\mathfrak{g}})_x$ 的中心 $Z(\hat{\mathfrak{g}}_x)$。

\begin{corollary}
    中心 $Z(\hat{\mathfrak{g}}_x)$ 同构于去心圆盘 $D_x^\times$ 上 ${}^LG$-opers 空间上函数代数 $\operatorname{Fun}\Op_{{}^LG}(D_x^\times)$。
\end{corollary}

我们还要记录下面这个有用的结果，其证明借自 [Beilinson and Drinfeld 1997] 注 3.7.11(iii)。注意群 $G((t))$ 自然作用在 $\widetilde{U}_{\kappa_c}(\hat{\mathfrak{g}})$ 上。

\begin{proposition}
    $G((t))$ 在 $Z(\hat{\mathfrak{g}})\subset\widetilde{U}_{\kappa_c}(\hat{\mathfrak{g}})$ 上的作用是平凡的。
\end{proposition}

\begin{proof}
设 $\widetilde{U}_{\kappa_c}(\hat{\mathfrak{g}})_{\le i}$ 是 $U_{\kappa_c}(\hat{\mathfrak{g}})$ 上 PBW 滤过在 $\widetilde{U}_{\kappa_c}(\hat{\mathfrak{g}})$ 上诱导的滤过的第 $i$ 项（见命题 4.3.4 的证明）。显然 $\widetilde{U}_{\kappa_c}(\hat{\mathfrak{g}})_{\le i}\,(i\ge0)$ 之并在 $\widetilde{U}_{\kappa_c}(\hat{\mathfrak{g}})$ 中稠密。此外，由中心作为代数 $\mathbb{C}[S_{i,[n]}]_{i=1,\dots,\ell;\,n\in\mathbb{Z}}$ 的完备化的描述可知，$Z(\hat{\mathfrak{g}})\cap\widetilde{U}_{\kappa_c}(\hat{\mathfrak{g}})_{\le i}\,(i\ge0)$ 之并在 $Z(\hat{\mathfrak{g}})$ 中稠密。因此只需证明 $G((t))$ 平凡地作用在 $Z(\hat{\mathfrak{g}})\cap\widetilde{U}_{\kappa_c}(\hat{\mathfrak{g}})$ 上（显然 $G((t))$ 保持每个 $\widetilde{U}_{\kappa_c}(\hat{\mathfrak{g}})_{\le i}$）。

在命题 4.3.4 中我们描述了 $Z(\hat{\mathfrak{g}})$ 关于上述滤过的分次伴随代数 $\operatorname{gr}Z(\hat{\mathfrak{g}})$，即 $\operatorname{gr}Z(\hat{\mathfrak{g}})$ 是某多项式代数的完备化。利用上面给出的该代数生成元 $\overline{P}_{i,n}$ 的显式公式，可知 $G((t))$ 平凡地作用在 $\operatorname{gr}Z(\hat{\mathfrak{g}})$ 上。因此 $G((t))$ 在 $Z(\hat{\mathfrak{g}})\cap\widetilde{U}_{\kappa_c}(\hat{\mathfrak{g}})$ 上的作用经由幂幺代数群 $U$ 分解。但其李代数 $\mathfrak{g}((t))$ 的相应作用经由 $U$ 的李代数的作用分解。由 $G$ 单连通的假设可知 $G((t))$ 连通，故 $U$ 也连通。

显然，连通幂幺李群的非平凡作用的微分也非平凡。因此，若 $G((t))$ 在 $Z(\hat{\mathfrak{g}})\cap\widetilde{U}_{\kappa_c}(\hat{\mathfrak{g}})$ 上的作用非平凡，则其李代数 $\mathfrak{g}((t))$ 的作用也非平凡。但 $\mathfrak{g}((t))$ 显然以 $0$ 作用在 $Z(\hat{\mathfrak{g}})$ 上。故 $G((t))$ 在 $Z(\hat{\mathfrak{g}})\cap\widetilde{U}_{\kappa_c}(\hat{\mathfrak{g}})$ 上、从而在 $Z(\hat{\mathfrak{g}})$ 上的作用平凡。证明完成。

注意，若不作 $G$ 单连通的假设，则 $G$ 的基本群是有限群 $\Gamma$，它等同于 $G((t))$ 的连通分支群。上述论证表明 $G((t))$ 的作用经由 $\Gamma$ 分解。但有限群没有非平凡的幂幺表示。因此即使不作 $G$ 单连通的假设，$G((t))$ 也平凡地作用在 $Z(\hat{\mathfrak{g}})$ 上。\transnote{卡点：命题 4.3.8 的证明为何先化到分次伴随代数。背景：$G((t))$ 保持 PBW 滤过；若它在 $\operatorname{gr}Z(\hat{\mathfrak{g}})$ 上作用平凡，则原作用落在某个幂幺群 $U$ 中，而连通幂幺群的非平凡作用其微分也非平凡。结论：但 $\mathfrak{g}((t))$ 显然以 $0$ 作用在中心上，故 $G((t))$ 的作用平凡；去掉单连通假设时作用至多经基本群 $\Gamma$ 分解，而有限群没有非平凡的幂幺表示。来源：原书 p.124 与 [Beilinson and Drinfeld 1997] 注 3.7.11(iii)（本书文献表）。}
\end{proof}

% [Tutor Note: 这个证明的关键归约是：$\mathrm{gr}Z$ 上作用平凡 $\Rightarrow$ 原作用落在幂幺群中。原文只说“factors through a unipotent algebraic group $U$”，其理由是对角可分化：若一个连通代数群的作用在分次伴随代数上平凡，则该作用可经一个滤过为幂幺的群分解，这一步请对照 [Beilinson and Drinfeld 1997] 注 3.7.11(iii)。]

\subsection{离开临界水平时的中心}

以上结果描述了临界水平 $\kappa=\kappa_c$ 下顶点代数 $V_\kappa(\mathfrak{g})$ 的中心与完备泛包络代数 $\widetilde{U}_\kappa(\hat{\mathfrak{g}})$ 的中心。离开临界水平又会怎样？

在命题 3.3.3 中我们已就顶点代数回答了此问题：中心平凡，即只含标量。现在为完整起见，对完备泛包络代数 $\widetilde{U}_\kappa(\hat{\mathfrak{g}})$ 回答此问题，并证明它同样平凡。因此就 $\widetilde{U}_\kappa(\hat{\mathfrak{g}})$ 的中心而言，临界水平之外我们没有遗漏任何东西。

\begin{proposition}
    当 $\kappa\ne\kappa_c$ 时，$\widetilde{U}_\kappa(\hat{\mathfrak{g}})$ 的中心由标量组成。
\end{proposition}

\begin{proof}
考虑李代数 $\hat{\mathfrak{g}}'_\kappa=\mathbb{C}d\ltimes\hat{\mathfrak{g}}_\kappa$，其中 $d$ 是这样的算子：它与 $\hat{\mathfrak{g}}_\kappa$ 的对易关系对应于向量场 $t\partial_t$ 的作用。设 $U'_\kappa(\hat{\mathfrak{g}})$ 是 $\hat{\mathfrak{g}}'_\kappa$ 的泛包络代数模去由 $\mathbf{1}-1$ 生成的理想所得的商，$\widetilde{U}'_\kappa(\hat{\mathfrak{g}})$ 是其按与 $\widetilde{U}_\kappa(\hat{\mathfrak{g}})$ 相同方式定义的完备化。

现在，代数 $\widetilde{U}_\kappa(\hat{\mathfrak{g}})$（$\kappa\ne\kappa_c$）含有算子 $L_0$，它与 $\hat{\mathfrak{g}}_\kappa$ 交换的方式如同 $-d$。因此 $L_0+d$ 是 $\widetilde{U}'_\kappa(\hat{\mathfrak{g}})$ 的中心元，且 $\widetilde{U}_\kappa(\hat{\mathfrak{g}})$ 同构于 $\widetilde{U}'_\kappa(\hat{\mathfrak{g}})$ 模去由 $L_0+d$ 生成的理想的商。

设 $Z_\kappa(\hat{\mathfrak{g}})$ 是 $\widetilde{U}_\kappa(\hat{\mathfrak{g}})$ 的中心，$Z'_\kappa(\hat{\mathfrak{g}})$ 是 $\widetilde{U}'_\kappa(\hat{\mathfrak{g}})$ 的中心。有自然嵌入 $\widetilde{U}_\kappa(\hat{\mathfrak{g}})\to\widetilde{U}'_\kappa(\hat{\mathfrak{g}})$。由于 $Z_\kappa(\hat{\mathfrak{g}})$ 的任一元素 $A$ 必满足 $[L_0,A]=0$，故 $A$ 在 $\widetilde{U}'_\kappa(\hat{\mathfrak{g}})$ 中的像也是中心元。于是
\[
    Z'_\kappa(\hat{\mathfrak{g}})=Z_\kappa(\hat{\mathfrak{g}})\otimes_{\mathbb{C}}\mathbb{C}[L_0+d].
\]

$Z'_\kappa(\hat{\mathfrak{g}})$ 的描述来自 V. Kac [1984] 的结果。按 [Kac 1984] 的推论 1，对 $\kappa\ne\kappa_c$，存在 $Z'_\kappa(\hat{\mathfrak{g}})$ 与超平面 $(\tilde{\mathfrak{h}})^*_\kappa$（对应于水平 $\kappa$，位于扩充 Cartan 子代数 $\tilde{\mathfrak{h}}=(\mathfrak{h}\otimes1)\oplus\mathbb{C}\mathbf{1}\oplus\mathbb{C}d$ 的对偶空间中）上关于仿射外尔群（经 $\rho$ 平移的）作用的不变函数代数之间的同构。但容易看出，满足此条件的 $(\tilde{\mathfrak{h}})^*_\kappa$ 上的多项式函数，由不变内积在 $\tilde{\mathfrak{h}}$ 上的限制所给出的多项式之幂次张成。这个二次多项式给出中心元 $L_0+d$。因此 $Z'_\kappa(\hat{\mathfrak{g}})=Z_\kappa(\hat{\mathfrak{g}})\otimes_{\mathbb{C}}\mathbb{C}[L_0+d]$，从而 $Z_\kappa(\hat{\mathfrak{g}})=\mathbb{C}$。\transnote{卡点：命题 4.3.9 的证明把中心化到 Kac 的仿射外尔群不变函数，其中“但容易看出”一句跳步较大。背景：$\hat{\mathfrak{g}}'_\kappa=\mathbb{C}d\ltimes\hat{\mathfrak{g}}_\kappa$ 的中心由 $L_0+d$ 与 $Z(\hat{\mathfrak{g}}_\kappa)$ 生成；Kac 1984 的推论 1 说 $Z'_\kappa(\hat{\mathfrak{g}})$ 同构于超平面 $(\tilde{\mathfrak{h}})^*_\kappa$ 上仿射外尔群（经 $\rho$ 平移）不变的多项式函数。结论：该超平面上不变多项式由不变内积限制所得二次多项式的幂次张成，故 $Z'_\kappa=Z_\kappa\otimes\mathbb{C}[L_0+d]$，从而 $Z_\kappa=\mathbb{C}$。来源：原书 p.125 与 [Kac 1984] 推论 1（本书文献表）。}
\end{proof}

% [Tutor Note: 这里的符号约定必须核对：原文说 $L_0$ 与 $\hat{\mathfrak{g}}_\kappa$ 交换“如同 $-d$”，于是 $L_0+d$ 才是中心元；若误读成“如同 $d$”，则应商去 $L_0-d$，结论不变但公式全变。另注意 $\mathbf{1}$ 是仿射代数的中心元（其像被商去：$\mathbf{1}-1$），与数字 $1$ 是两回事。]
